\documentclass[11pt]{article}

\usepackage[spanish,es-tabla]{babel}
\usepackage[utf8]{inputenc}
\usepackage[T1]{fontenc}
\usepackage{amsmath,amssymb}
\usepackage{booktabs}
\usepackage{array}
\usepackage{tabularx}
\usepackage{longtable}
\usepackage{geometry}
\usepackage{enumitem}
\usepackage[most]{tcolorbox}
\usepackage{hyperref}

\geometry{letterpaper, margin=1.85cm}
\setlength{\parindent}{0pt}
\setlength{\parskip}{0.55em}
\setlist[itemize]{leftmargin=1.35em,itemsep=0.25em}
\setlist[enumerate]{leftmargin=1.45em,itemsep=0.25em}

\hypersetup{colorlinks=true, linkcolor=blue!50!black, urlcolor=blue!50!black}

\tcbset{
  colback=white,
  colframe=black!60,
  boxrule=0.45pt,
  arc=1.2mm,
  left=1.5mm,
  right=1.5mm,
  top=1mm,
  bottom=1mm,
  breakable
}

\newtcolorbox{pregunta}[1]{title=\textbf{Pregunta anterior: #1}, colback=green!3, colframe=green!45!black}
\newtcolorbox{enunciado}[1]{title=\textbf{Enunciado completo: #1}, colback=green!3, colframe=green!45!black}
\newtcolorbox{materia}[1]{title=\textbf{Materia que activa: #1}, colback=blue!3, colframe=blue!50!black}
\newtcolorbox{pauta}[1]{title=\textbf{Pauta / receta: #1}, colback=gray!5, colframe=black!60}
\newtcolorbox{pautacompleta}[1]{title=\textbf{Pauta completa paso a paso: #1}, colback=gray!5, colframe=black!60}
\newtcolorbox{alerta}[1]{title=\textbf{Ojo de examen: #1}, colback=red!3, colframe=red!55!black}

\newcommand{\E}{\mathbb{E}}
\newcommand{\Var}{\operatorname{Var}}
\newcommand{\Prob}{\mathbb{P}}

\title{\textbf{Preguntas de Examenes Anteriores por Materia}}
\author{Gestion de Operaciones}
\date{}

\begin{document}
\maketitle
\tableofcontents
\newpage

\section{Mapa General}

El objetivo de este compilado es revisar preguntas pasadas por eje tematico, no por año. Para cada bloque se conserva el tipo de enunciado, la materia que se debe reconocer y una pauta convertida en receta de resolucion. Los ejes siguen el temario del examen: bodegas y centros de distribucion, calidad, coordinacion en la cadena, TPS/JIT/Lean, casos Toyota/Zara y arrastre de I1/I2.

\begin{center}
\begin{tabularx}{\linewidth}{p{0.18\linewidth}X p{0.20\linewidth}}
\toprule
Eje & Preguntas anteriores mas utiles & Prioridad \\
\midrule
Calidad & 2015 tarjetas JIT + control, 2016 vasos/proveedores, 2020 pan, 2021 $\bar X/R$, 2024 botellas & Muy alta \\
Bodegas/CD & 2016 picking frontal y volumen optimo, 2021 tamaño de bodega y picking rapido, 2022 profundidad de pallets, 2023 tamaño de CD & Muy alta \\
Coordinacion & 2017 demostracion $Q_C=2Q_D$, 2020 Burger Mac, 2022 franquicias, 2023 cafeteria & Muy alta \\
TPS/JIT/Lean & 2015 La Meta + CB, 2015/2018/2019 JIT/Lean, 2018 Six Sigma/TQM, 2024 proceso con desecho & Alta \\
Variabilidad/colas & 2018 autopista G/G/c, 2019 VUT proceso, 2020 CB con variabilidad, 2021 M/M/1, 2023 G/G/1 CD, 2024 fallas & Alta \\
Arrastre I1/I2 & Pronosticos, EOQ, WW, PERT, MRP y optimizacion & Media-alta \\
\bottomrule
\end{tabularx}
\end{center}

\begin{alerta}{uso durante el examen}
Si aparece una pregunta parecida, no busques el año: busca el eje. La misma matematica aparece con nombres distintos: cafeterias, hamburguesas, bodegas, botellas o tarjetas.
\end{alerta}

\section{Calidad}

\subsection{Control estadistico de procesos: pan, tarjetas y botellas}

\begin{materia}{graficos de control}
Cuando el enunciado entrega subgrupos con promedio y recorrido, define para cada subgrupo $k$:
\[
  \bar x_k=\frac{1}{n}\sum_{j=1}^n x_{kj}, \qquad
  R_k=\max_j x_{kj}-\min_j x_{kj}.
\]
Luego:
\[
  \bar{\bar x}=\frac{1}{m}\sum_{k=1}^{m}\bar x_k,
  \qquad
  \bar R=\frac{1}{m}\sum_{k=1}^{m}R_k.
\]
Con constantes de tabla para tamaño de muestra $n$:
\[
  LCS_{\bar X}=\bar{\bar x}+A_2\bar R,\quad
  LCI_{\bar X}=\bar{\bar x}-A_2\bar R,
\]
\[
  LCS_R=D_4\bar R,\quad
  LCI_R=D_3\bar R.
\]
Si el enunciado pide un nivel de confianza usando agregados de medias, usa:
\[
  LCS=\bar x+z s,\qquad LCI=\bar x-zs,
\]
donde $s^2=\frac{1}{m}\sum \bar x_k^2-\bar x^2$.
\end{materia}

\begin{enunciado}{Examen 2020, pan}
Usted decide controlar el proceso de produccion de pan y decide hacer un esquema de muestra de 30 panes a los cuales les mide el peso en gramos. Los resultados de los 27 grupos de muestra se detallan a continuacion, con el respectivo promedio de cada muestra y recorrido.

\begin{center}
\small
\begin{tabular}{c|cc@{\hspace{1.0em}}c|cc@{\hspace{1.0em}}c|cc}
\toprule
\# & Promedio & Recorrido & \# & Promedio & Recorrido & \# & Promedio & Recorrido\\
\midrule
1 & 8.67 & 1.6 & 10 & 9.52 & 1.1 & 19 & 9.61 & 1.2\\
2 & 9.70 & 1.5 & 11 & 10.61 & 1.2 & 20 & 11.65 & 1.6\\
3 & 9.73 & 2.0 & 12 & 11.56 & 0.6 & 21 & 11.00 & 0.9\\
4 & 10.00 & 0.1 & 13 & 10.10 & 1.1 & 22 & 10.72 & 1.5\\
5 & 9.55 & 1.3 & 14 & 9.46 & 0.3 & 23 & 9.41 & 1.8\\
6 & 10.67 & 1.5 & 15 & 10.18 & 0.0 & 24 & 8.67 & 1.3\\
7 & 10.90 & 0.0 & 16 & 11.26 & 0.6 & 25 & 12.30 & 0.0\\
8 & 9.91 & 1.8 & 17 & 8.03 & 0.4 & 26 & 9.28 & 1.2\\
9 & 9.91 & 1.8 & 18 & 10.81 & 0.6 & 27 & 9.93 & 0.9\\
\bottomrule
\end{tabular}
\end{center}

Ademas se sabe que:
\[
  \sum \text{promedios}=273.14,\qquad
  \sum \text{recorrido}=27.9,
\]
\[
  \sum(\text{promedios})^2=2788.03,\qquad
  \sum(\text{recorrido})^2=38.67.
\]
Hint:
\[
  \sigma^2=\frac{1}{n}\sum X^2-\bar X^2=E[X^2]-E[X]^2.
\]

\textbf{Pregunta.} Con esta informacion desarrolle los graficos de control del proceso para un control $96\%$.
\end{enunciado}

\begin{pautacompleta}{Examen 2020, pan}
La pauta usa normalidad sobre los promedios de los 27 grupos:
\[
  \bar x=\frac{273.14}{27}=10.116.
\]
La varianza de los promedios se calcula como:
\[
  s^2=\frac{2788.03}{27}-(10.116)^2\approx 0.921.
\]
Para $96\%$ bilateral, la pauta usa $z\simeq 2.05$. Entonces:
\[
  LCS=10.116+2.05\sqrt{0.921},\qquad
  LCI=10.116-2.05\sqrt{0.921}.
\]
Se eliminan las muestras fuera de limite, en la pauta las muestras 17 y 25, y se recalcula:
\[
  \sum \bar x_k=273.14-8.03-12.30=251.81,
\]
\[
  \sum \bar x_k^2=2788.03-8.03^2-12.30^2=2572.26.
\]
Con $m=25$:
\[
  \bar x_{\text{nuevo}}=\frac{251.81}{25}=10.112,\qquad
  s^2_{\text{nuevo}}=\frac{2572.26}{25}-(10.112)^2\approx 0.793.
\]
Los limites finales son:
\[
  LCS=10.112+2.05\sqrt{0.793},\qquad
  LCI=10.112-2.05\sqrt{0.793}.
\]
Todo lo restante queda dentro de rango, por lo que esos son los limites finales del proceso segun la pauta.
\end{pautacompleta}

\begin{enunciado}{Examen 2024, botellas de vidrio}
Usted decide controlar un proceso de produccion de botellas de vidrio y decide hacer un esquema de muestras de 20 botellas cada una a las cuales les mide el grosor en milimetros. Los resultados de 25 grupos de muestra se detallan a continuacion, con su respectivo promedio, valor minimo y maximo de cada muestra.

\begin{center}
\small
\begin{tabular}{c|ccc@{\hspace{1.2em}}c|ccc}
\toprule
Muestra & Promedio & Minimo & Maximo & Muestra & Promedio & Minimo & Maximo\\
\midrule
1 & 25.1 & 24.0 & 26.2 & 14 & 26.1 & 24.9 & 27.3\\
2 & 25.8 & 24.5 & 27.1 & 15 & 24.6 & 23.4 & 25.8\\
3 & 24.7 & 23.9 & 25.5 & 16 & 27.5 & 26.2 & 28.8\\
4 & 26.4 & 25.2 & 27.8 & 17 & 25.7 & 24.5 & 26.9\\
5 & 25.3 & 24.3 & 26.3 & 18 & 24.8 & 23.6 & 26.0\\
6 & 24.5 & 23.5 & 25.5 & 19 & 26.3 & 25.0 & 27.6\\
7 & 25.9 & 24.6 & 27.2 & 20 & 25.0 & 23.9 & 26.1\\
8 & 25.2 & 24.2 & 26.2 & 21 & 24.5 & 23.3 & 25.7\\
9 & 24.9 & 23.7 & 26.1 & 22 & 27.6 & 26.3 & 28.9\\
10 & 26.0 & 24.8 & 27.2 & 23 & 25.6 & 24.4 & 26.8\\
11 & 26.4 & 25.1 & 27.7 & 24 & 25.1 & 23.8 & 26.4\\
12 & 25.4 & 24.4 & 26.4 & 25 & 24.4 & 23.1 & 25.7\\
13 & 24.2 & 23.0 & 25.4 & & & &\\
\bottomrule
\end{tabular}
\end{center}

Datos agregados de la tabla:
\[
\sum \text{Promedio}=637.1,\quad
\sum \text{Minimo}=607.6,\quad
\sum \text{Maximo}=666.6,
\]
\[
\sum(\text{Promedio})^2=16255.53,\quad
\sum(\text{Minimo})^2=14784.56,\quad
\sum(\text{Maximo})^2=17796.96.
\]

\begin{enumerate}[label=\alph*)]
  \item Con esta informacion determine los limites de control del grosor de botella.
  \item Suponga ahora que calcula una nueva muestra, con media $25.3$ mm y recorrido $0.8$ mm. ¿Que puede decir del proceso? ¿Se encuentra bajo control? ¿Por que?
  \item Si ahora le informan que las muestras tomadas son de 30 botellas cada una, con la misma informacion de la tabla inicial. En base a esta nueva informacion, desarrolle los graficos de control del proceso para un control $90\%$. ¿Como cambia lo obtenido en a)?
  \item Para el tamaño de muestra de 30 botellas, ¿cual es el porcentaje de confianza minimo necesario para que no se deban eliminar muestras en el desarrollo de graficos del inciso anterior?
  \item ¿Este porcentaje es alto o bajo? ¿Que indica esto sobre el proceso de produccion? ¿Que efecto tiene eliminar muestras con respecto al porcentaje de confianza en este proceso?
\end{enumerate}
\end{enunciado}

\begin{pautacompleta}{Examen 2024, botellas de vidrio}
Para $n=20$, la pauta usa constantes:
\[
  A_2=0.18,\qquad D_3=0.415,\qquad D_4=1.585.
\]
El promedio global de los promedios es:
\[
  \bar{\bar x}=\frac{637.1}{25}=25.484.
\]
El recorrido de cada muestra es $R_k=\max_k-\min_k$. Con los datos agregados:
\[
  \bar R=\frac{\sum \max-\sum \min}{25}
  =
  \frac{666.6-607.6}{25}
  =
  2.36.
\]
Los limites para el grafico de promedios son:
\[
  LCI_{\bar X}=25.484-0.18(2.36)=25.06,\qquad
  LCS_{\bar X}=25.484+0.18(2.36)=25.90,
\]
y los limites para el grafico de recorridos:
\[
  LCI_R=0.415(2.36)=0.979,\qquad
  LCS_R=1.585(2.36)=3.741.
\]

La nueva muestra tiene media $25.3$ y recorrido $0.8$. La media cae dentro de $[25.06,25.90]$, pero el recorrido queda bajo el limite inferior:
\[
  0.8<0.979.
\]
Por lo tanto, el proceso no se encuentra bajo control segun el grafico de recorrido.

Para $n=30$ y confianza $90\%$, la pauta usa $z=1.64$ sobre la desviacion de los promedios:
\[
  s=\sqrt{\frac{\sum \bar x_k^2}{25}-(25.484)^2}=0.887.
\]
Limites iniciales:
\[
  25.484\pm 1.64(0.887)=[24.046,26.922].
\]
Se eliminan las muestras 16 y 22, se recalcula $\bar x=25.304$, $s=0.685$, y quedan:
\[
  25.304\pm 1.64(0.685)=[24.181,26.427].
\]
Para que no se eliminen muestras, se iguala el limite a los promedios extremos. El promedio inferior extremo es $24.2$ y el superior extremo es $27.6$:
\[
  Z_{\inf}=\frac{25.484-24.2}{0.887}=1.44,\qquad
  Z_{\sup}=\frac{27.6-25.484}{0.887}=2.39.
\]
El valor que protege ambos extremos es $Z=2.39$, por lo que la confianza minima aproximada es:
\[
  2\Phi(2.39)-1\simeq 98.32\%.
\]
Ese porcentaje es alto: para no eliminar muestras se necesita un intervalo muy ancho, lo que sugiere un proceso con variabilidad relevante. Eliminar muestras estrecha los limites y aumenta el riesgo de rechazar muestras que antes se aceptaban; en lenguaje de muestreo, aumenta el riesgo del productor.
\end{pautacompleta}

\begin{enunciado}{Examen 2015, Tarjetas ABC}
La empresa Tarjetas ABC desea establecer un plan de produccion JIT. La demanda diaria registrada es de 200 tarjetas telefonicas por hora. El proceso de produccion de estas tarjetas pasa por 3 grandes operaciones antes del control de calidad ubicado al final de la linea: impresion de las leyendas (P1), inclusion del chip (P2) y cortado de la tarjeta (P3).

La empresa cuenta con un registro de los tiempos promedios de procesamiento ($t_{pi}$) por operacion, ademas de los tiempos de envio de los Kanbans ($t_{ki}$) y tiempos de envio de los lotes ($t_{vi}$).

\begin{center}
\begin{tabular}{c|cccc}
\toprule
Operacion & Lote $C$ & $t_{pi}$ seg & $t_{ki}$ seg & $t_{vi}$ seg\\
\midrule
P1 & 200 & 85 & 45 & 200\\
P2 & 250 & 78 & 67 & 300\\
P3 & 300 & 50 & 92 & 150\\
\bottomrule
\end{tabular}
\end{center}

Los registros historicos del control de calidad indican que en promedio un $15\%$ de las unidades son descartadas.

\begin{enumerate}[label=\alph*)]
  \item A partir de los datos anteriores, calcule el numero de Kanbans necesarios en el proceso productivo.
\end{enumerate}

Las investigaciones de la empresa han determinado que el proceso P3 es el que esta actualmente generando el $15\%$ del descarte de las tarjetas, las cuales no estan saliendo con los tamaños adecuados. Se realizo un muestreo del largo de las tarjetas de 5 lotes recibidos durante los ultimos 5 dias.

\begin{center}
\begin{tabular}{c|ccccc}
\toprule
Dia & \multicolumn{5}{c}{Largo (milimetros)}\\
\midrule
1 & 70 & 56 & 49 & 67 & 61\\
2 & 65 & 47 & 70 & 70 & 68\\
3 & 70 & 49 & 42 & 68 & 54\\
4 & 50 & 47 & 52 & 67 & 50\\
5 & 48 & 65 & 51 & 50 & 65\\
\bottomrule
\end{tabular}
\end{center}

A partir de lo anterior:
\begin{enumerate}[label=\alph*),resume]
  \item Calcule los limites de control, eliminando los outliers. Realice los graficos correspondientes.
  \item Si el dia 6 usted toma una muestra con los siguientes resultados: $63$, $54$, $43$, $69$ y $65$. ¿Que puede decir de la muestra? ¿Esta el proceso bajo control?
  \item El implementar el control reduce las fallas del sistema a solo un $5\%$. Con esta informacion, ¿cambia el numero Kanbans? Argumente.
\end{enumerate}
\end{enunciado}

\begin{pautacompleta}{Examen 2015, Tarjetas ABC}
Primero se calcula el lead time total de reposicion por proceso:
\[
  L_i=t_{p,i}+t_{k,i}+t_{v,i}.
\]
Por lo tanto:
\[
  L_1=85+45+200=330,\qquad
  L_2=78+67+300=445,\qquad
  L_3=50+92+150=292.
\]
Con descarte de $15\%$, el rendimiento es $y=0.85$. La pauta calcula:
\[
  N_i=\frac{D}{0.85}\frac{L_i}{C_i}.
\]
Resultados reportados:
\[
  N_1=389,\qquad N_2=419,\qquad N_3=230.
\]

Para calidad, calcula $\bar x_k$ y $R_k$ por dia:
\[
\begin{array}{c|cc}
\text{Dia} & \bar x_k & R_k\\
\hline
1&60.6&21\\
2&64.0&23\\
3&54.6&28\\
4&53.2&20\\
5&55.8&17\\
\hline
\text{Promedio}&58.04&21.8
\end{array}
\]
Con $A_2=0.58$, $D_3=0$, $D_4=2.11$:
\[
  LCS_{\bar X}=58.04+0.58(21.8),\qquad
  LCI_{\bar X}=58.04-0.58(21.8),
\]
\[
  LCS_R=2.11(21.8),\qquad LCI_R=0.
\]
Despues de eliminar el dato fuera de control, la pauta usa $\bar{\bar x}=58.8$ y $\bar R=24.6$. La muestra del dia 6 tiene:
\[
  \bar x_6=\frac{63+54+43+69+65}{5}=58.8,\qquad R_6=69-43=26.
\]
Se compara contra los limites recalculados. La regla de examen es: media dentro y recorrido dentro implica muestra bajo control; si uno sale, no.

Si las fallas bajan a $5\%$, el nuevo rendimiento es $y=0.95$. Como lo demas queda constante:
\[
  N_{i,\text{nuevo}}=N_{i,\text{viejo}}\frac{0.85}{0.95}.
\]
La pauta reporta:
\[
  N_1=348,\qquad N_2=375,\qquad N_3=205.
\]
\end{pautacompleta}

\subsection{Muestreo de aceptacion, AQL y LPTD}

\begin{materia}{curva OC y riesgos}
En muestreo de aceptacion simple, se define $X$ como numero de defectuosos en una muestra de tamaño $n$:
\[
  X\sim Bin(n,p),\qquad
  P_{\text{aceptar}}(p)=\Prob(X\le c)=\sum_{x=0}^{c}\binom{n}{x}p^x(1-p)^{n-x}.
\]
La curva OC es $P_{\text{aceptar}}(p)$ como funcion de $p$. En AQL se quiere aceptar con probabilidad alta; en LPTD se quiere aceptar con probabilidad baja. El riesgo del productor $\alpha$ es rechazar un lote bueno; el riesgo del consumidor $\beta$ es aceptar un lote malo.
\end{materia}

\begin{enunciado}{Examen 2019, verdadero/falso}
Seccion verdadero o falso. Indique si la siguiente afirmacion es verdadera o falsa. En caso de ser falsa, indique la razon:

``Al disminuir el numero o nivel maximo de aceptacion en un muestreo, por ejemplo de 5 unidades a 3 unidades, se aumenta el riesgo del productor o error tipo $\alpha$.''
\end{enunciado}

\begin{pautacompleta}{pauta 2019}
Verdadero. Al bajar $c$, el plan es mas estricto. Eso reduce $P_{\text{aceptar}}(p)$ para todo $p$, incluyendo lotes buenos en AQL. Por lo tanto aumenta la probabilidad de rechazar un lote bueno:
\[
  \alpha=1-P_{\text{aceptar}}(p_{AQL}).
\]
\end{pautacompleta}

\subsection{Calidad de proveedores y mejora de proceso}

\begin{enunciado}{Examen 2016, vasos termicos y proveedor}
Usted es el titular de una empresa de cafe preparado llamada ``Juana Valdez Cafe''. Habitualmente compra el grano, y en cada uno de sus tres locales lo muele y sirve en vasos de carton termico. Debido a las quejas de los clientes, han registrado que los vasos no siempre se comportan termicamente bien. Por ello ha decidido hacer un estudio de proveedores.

Usted ha determinado que los vasos deberian tolerar un minimo de $90^\circ\text{C}$ para que sean seguros y coincida con lo que exige a sus proveedores. Por otro lado, tiene un historico de muestreos de recepcion en donde, en promedio, los vasos tienen una resistencia de $95^\circ\text{C}$, con una desviacion estandar de $2^\circ\text{C}$. Actualmente su proveedor ``Kartoon'' le entrega en lotes de 5000 vasos.

Usted ha decidido iniciar un estudio del problema y esta analizando como definir bien el problema para poder llegar a una solucion. Revisando sus apuntes de Gestion de Operaciones se encuentra con 4 preguntas que le pueden ayudar en su trabajo. A continuacion, desarrolle las respuestas para cada pregunta:

\begin{enumerate}[label=\alph*)]
  \item Determine la informacion que necesitaria recopilar del proceso y del proveedor para analizar el problema.
  \item Establezca que metodologia utilizaria para mejorar la calidad del proceso y como evaluaria esta.
  \item Realice el mismo analisis, pero ahora enfocado en la recepcion de los insumos del proveedor.
  \item La aplicacion de ambos cambios llevaria a mejorar la calidad del producto final?
\end{enumerate}
\end{enunciado}

\begin{pautacompleta}{pauta completa 2016}
Esta pregunta no es de calculo directo; es de diagnostico de calidad. La respuesta debe separar proceso interno, proveedor, sistema de medicion, metodologia de mejora y recepcion de lotes.

Para el inciso a), en el proceso interno hay que registrar variables que permitan observar el problema y sus posibles causas:
\begin{itemize}
  \item Temperatura de llenado de los vasos, separada por tipo de bebida.
  \item Temperaturas indicadas por los equipos en sus visores.
  \item Lote de producto ingresado a linea, usado para despachos de bebidas calientes.
  \item Volumen de llenado habitual, porque el problema podria venir de salpicaduras o sobrellenado y no solamente de resistencia termica del vaso.
  \item Registro de quejas de clientes con detalle: local, fecha, bebida, temperatura, lote del vaso y descripcion del incidente.
\end{itemize}

Respecto del proveedor, hay que recuperar informacion del comportamiento historico de su proceso:
\begin{itemize}
  \item Ensayos de resistencia termica en produccion.
  \item Registros historicos de control estadistico de proceso, para ver si existen curvas fuera de control o lotes cercanos a especificacion.
  \item Tasa media de defectos.
  \item Tamaño de lote, que en este caso es 5000 vasos.
  \item Especificaciones, tolerancias aceptadas y criterios usados para liberar lotes.
  \item Productos mas criticos y trazabilidad de lote.
\end{itemize}

Tambien se debe homologar el sistema de medicion. No basta con comparar la media historica $95^\circ\text{C}$ contra el minimo $90^\circ\text{C}$ si el proveedor y la cafeteria miden de forma distinta. La pauta exige cotejar la practica de muestreo usada por el proveedor para liberar producto con el esquema de aceptacion de la cafeteria. Esto corresponde a revisar practicas analiticas y, si corresponde, hacer un estudio R\&R.

Para el inciso b), una metodologia natural es DMAIC:
\[
  \text{Definir} \rightarrow \text{Medir} \rightarrow \text{Analizar} \rightarrow \text{Mejorar} \rightarrow \text{Controlar}.
\]
Aunque no se implemente Six Sigma completo, DMAIC entrega una estructura valida. En este caso:
\begin{itemize}
  \item Definir el defecto: vaso que no tolera el uso termico esperado o genera queja de cliente.
  \item Medir temperaturas de llenado, resistencia del vaso, lotes, maquinas y locales.
  \item Analizar si la causa esta en el proveedor, en el diseño del vaso, en la variabilidad de temperatura, en el volumen de llenado o en la manipulacion.
  \item Mejorar mediante estandarizacion operacional, desarrollo de proveedores y ajustes de proceso.
  \item Controlar con protocolos de medicion, auditorias periodicas y seguimiento SPC.
\end{itemize}

La pauta menciona tres lineas de mejora: reducir dispersion en la temperatura de llenado, desarrollar proveedores para que la resistencia termica venga desde el diseño del vaso, y estandarizar operaciones dentro de la cafeteria con controles visuales.

La evaluacion de estas medidas debe hacerse con un plan de analisis que pruebe la funcionalidad del vaso en condiciones extremas. Con esos datos se puede construir un historico tipo SPC. Ademas, deben realizarse auditorias periodicas al proveedor, revisando fabricacion en linea y registros historicos.

Para el inciso c), enfocado en recepcion de insumos, se plantea un sistema de muestreo de aceptacion. El proceso mecanico es:
\begin{enumerate}
  \item Elegir el tipo de muestreo: simple, doble o triple.
  \item Usar el tamaño de lote $N=5000$, el nivel de calidad aceptable y la diferencia que se desea detectar para determinar el tamaño de muestra $n$.
  \item Definir el criterio de aceptacion $c$, es decir, el maximo numero de vasos defectuosos permitido en la muestra para aceptar el lote.
  \item Definir como se hara el pickeo de muestras segun la constitucion de pallets, evitando tomar muestras sesgadas de una sola zona del lote.
  \item Aplicar los ensayos de resistencia termica y decidir aceptar o rechazar el lote.
\end{enumerate}

Matematicamente, si $X$ es el numero de vasos defectuosos en la muestra:
\[
  X\sim Bin(n,p),
  \qquad
  P(\text{aceptar lote})=\Prob(X\le c).
\]
La especificacion relevante es:
\[
  \text{resistencia termica}\geq 90^\circ\text{C}.
\]
Un vaso bajo $90^\circ\text{C}$ se considera no conforme para este criterio.

Para el inciso d), la respuesta correcta no es afirmar mecanicamente que ambos cambios siempre mejoran. La pauta indica que el producto final deberia mejorar, pero la mejora podria no ser sostenible e incluso tener efectos negativos si solo se endurece la recepcion. Por ejemplo, si aumentan los rechazos sin mejorar la calidad de origen, podria haber quiebres de stock.

La mejora real exige dos condiciones:
\begin{itemize}
  \item Desarrollar con el proveedor una alternativa que produzca calidad desde el proceso de fabricacion, no solo un control mas estricto de tipo pasa/no pasa.
  \item Controlar suficientemente el procedimiento de la cafeteria para evitar que el uso real lleve los vasos al borde de la especificacion. Los limites del proceso interno deben quedar dentro de las especificaciones de uso del vaso.
\end{itemize}

La frase clave para responder en examen es: el control de recepcion detecta lotes malos, pero el desarrollo del proveedor y el control del proceso interno reducen la probabilidad de producir o usar mal el vaso.
\end{pautacompleta}

\begin{alerta}{frase para examen}
En calidad de proveedores, distinguir entre \emph{control de recepcion} y \emph{desarrollo de proveedor}. El primero detecta defectos; el segundo reduce la probabilidad de que aparezcan.
\end{alerta}

\section{Bodegas y Centros de Distribucion}

\subsection{Picking rapido y asignacion de espacio}

\begin{materia}{volumen optimo para picking rapido}
Sea $I$ el conjunto de SKU. Para cada SKU $i\in I$:
\begin{itemize}
  \item $f_i$ es el flujo volumetrico mensual en $m^3/\text{mes}$.
  \item $v_i$ es el volumen asignado al SKU $i$.
  \item $V$ es el volumen total disponible.
  \item $c_r$ es el costo por reposicion.
\end{itemize}
El modelo clasico es:
\[
  \min_{v_i\ge 0}\; c_r\sum_{i\in I}\frac{f_i}{v_i}
  \quad \text{s.a.}\quad
  \sum_{i\in I}v_i=V.
\]
La condicion de optimalidad entrega:
\[
  v_i^\star=V\frac{\sqrt{f_i}}{\sum_{j\in I}\sqrt{f_j}}.
\]
Si los costos de reposicion son distintos:
\[
  \min \sum_i c_{r,i}\frac{f_i}{v_i},
  \qquad
  v_i^\star=V\frac{\sqrt{c_{r,i}f_i}}{\sum_j\sqrt{c_{r,j}f_j}}.
\]
\end{materia}

\begin{enunciado}{Examen 2016, centro de distribucion}
En otro centro de distribucion se tiene un area de Picking rapido cuyo espacio total disponible es de $1500\ \text{m}^3$, en la cual se busca disponer tres SKU: X, Y, Z y W.

\begin{center}
\begin{tabular}{c|ccc}
\toprule
SKU & UNIDADES/MES & UNIDADES/CAJA & M3/CAJA\\
\midrule
X & 3000 & 20 & 3\\
Y & 3600 & 24 & 2\\
Z & 2040 & 36 & 1.5\\
W & 1000 & 50 & 5\\
\bottomrule
\end{tabular}
\end{center}

\begin{enumerate}[label=\alph*)]
  \item Plantee el problema de optimizacion que minimiza el costo total del area de picking rapido. Suponga que el costo de \emph{restock} es $c_r$.
  \item Con los datos de la tabla anterior calcule el espacio optimo y la frecuencia de \emph{restock} para cada SKU.
  \item Calcule el espacio asignado y el \emph{restock} para cada SKU considerando las siguientes configuraciones: igual volumen para los SKU y con igual frecuencia de \emph{restock}.
\end{enumerate}
\end{enunciado}

\begin{pautacompleta}{pauta paso a paso 2016}
Primero convierte cada SKU a flujo volumetrico:
\[
  f_i=\frac{\text{unidades/mes}_i}{\text{unidades/caja}_i}\cdot m_i^3/\text{caja}.
\]
Con los datos de la tabla:
\[
  f_X=\frac{3000}{20}\cdot 3=450,\qquad
  f_Y=\frac{3600}{24}\cdot 2=300,
\]
\[
  f_Z=\frac{2040}{36}\cdot 1.5=85,\qquad
  f_W=\frac{1000}{50}\cdot 5=100.
\]
El modelo que minimiza el costo total de reposiciones es:
\[
  \min_{v_i\geq 0}\; c_r\sum_{i\in I}\frac{f_i}{v_i}
  \quad\text{s.a.}\quad
  \sum_{i\in I} v_i=1500,
\]
donde $I=\{X,Y,Z,W\}$, $f_i$ es el flujo volumetrico mensual del SKU $i$ y $v_i$ es el volumen asignado al SKU $i$.

La condicion de optimalidad entrega:
\[
  v_i^\star=1500\frac{\sqrt{f_i}}{\sqrt{450}+\sqrt{300}+\sqrt{85}+\sqrt{100}}.
\]
La suma de raices es:
\[
  \sqrt{450}+\sqrt{300}+\sqrt{85}+\sqrt{100}=57.754.
\]
Los volumenes optimos y frecuencias son:
\[
\begin{array}{c|cc}
\text{SKU} & v_i^\star & \text{restock } f_i/v_i\\
\hline
X&550.96&0.82\\
Y&449.86&0.67\\
Z&239.46&0.35\\
W&259.73&0.39
\end{array}
\]

Para igual volumen:
\[
  v_i=\frac{1500}{4}=375\ \text{m}^3\quad \forall i.
\]
Las frecuencias de reposicion son:
\[
\begin{array}{c|cc}
\text{SKU} & v_i & f_i/v_i\\
\hline
X&375&1.20\\
Y&375&0.80\\
Z&375&0.23\\
W&375&0.27
\end{array}
\]

Para igual frecuencia, se exige $f_i/v_i=k$ para todo $i$, por lo tanto:
\[
  v_i=V\frac{f_i}{\sum_j f_j}.
\]
Como $\sum_j f_j=935$, se obtiene:
\[
\begin{array}{c|cc}
\text{SKU} & v_i & f_i/v_i\\
\hline
X&1500(450/935)=721.93&0.623\\
Y&1500(300/935)=481.28&0.623\\
Z&1500(85/935)=136.36&0.623\\
W&1500(100/935)=160.43&0.623
\end{array}
\]
\end{pautacompleta}

\begin{enunciado}{Examen 2021, picking rapido}
Su empresa ha construido una zona de \emph{picking} rapido de $20\ \text{m}^3$ y se le pide a usted ubicar los SKU que se detallan en la tabla a continuacion.

\begin{center}
\begin{tabular}{c|cc}
\toprule
SKU & cajas/mes & $\text{m}^3$/caja\\
\midrule
A & 40 & 2\\
B & 20 & 1.5\\
C & 45 & 1\\
\bottomrule
\end{tabular}
\end{center}

Si el costo de reposicion es de \$3 por reposicion para cada SKU y usando una politica de volumen optimo y otra de igual tiempo, determine:

\begin{enumerate}[label=\alph*)]
  \item Para cada politica, cuanto espacio asigna a cada uno y cual es el costo de cada politica.
  \item Si los costos de cada reposicion son distintos para cada SKU, siendo \$3 para SKU A, \$2 para SKU B y \$3 para SKU C, plantee un modelo que permita determinar la asignacion optima para cada SKU.
\end{enumerate}
\end{enunciado}

\begin{pautacompleta}{pauta paso a paso 2021}
Se calculan los flujos:
\[
  f_A=40(2)=80,\qquad f_B=20(1.5)=30,\qquad f_C=45(1)=45.
\]
Con volumen optimo:
\[
  v_i^\star=V\frac{\sqrt{f_i}}{\sum_j \sqrt{f_j}}.
\]
Con $V=20$:
\[
  \sum_j\sqrt{f_j}=\sqrt{80}+\sqrt{30}+\sqrt{45}=21.131.
\]
Por lo tanto:
\[
  v_A^\star=20\frac{\sqrt{80}}{21.131}=8.47,\quad
  v_B^\star=20\frac{\sqrt{30}}{21.131}=5.18,\quad
  v_C^\star=20\frac{\sqrt{45}}{21.131}=6.35.
\]
Las reposiciones mensuales son:
\[
  r_A=\frac{80}{8.47}=9.45,\quad
  r_B=\frac{30}{5.18}=5.79,\quad
  r_C=\frac{45}{6.35}=7.09.
\]
El costo mensual de reposicion con volumen optimo es:
\[
  CT_{\text{opt}}=3(r_A+r_B+r_C)=3(22.33)=66.99.
\]

En igual tiempo se igualan las frecuencias:
\[
  \frac{f_A}{v_A}=\frac{f_B}{v_B}=\frac{f_C}{v_C}.
\]
Entonces se asigna proporcional a flujo:
\[
  v_i=V\frac{f_i}{\sum_j f_j},\qquad \sum_j f_j=155.
\]
La tabla queda:
\[
\begin{array}{c|ccc}
\text{SKU} & f_i & v_i & f_i/v_i\\
\hline
A&80&20(80/155)=10.32&7.75\\
B&30&20(30/155)=3.87&7.75\\
C&45&20(45/155)=5.81&7.75
\end{array}
\]
El costo mensual es:
\[
  CT_{\text{igual tiempo}}=3(7.75+7.75+7.75)=69.75.
\]

\textbf{Nota de consistencia.} La pauta escrita usa $V=25$ en los calculos numericos y obtiene $v_A=10.58$, $v_B=6.48$, $v_C=7.94$, costo $53.575$, y para igual tiempo volumenes $12.90$, $4.84$, $7.26$ con costo $55.8$. El enunciado transcrito dice $20\ \text{m}^3$; por eso aqui se recalcula con $V=20$. La proporcion y el algoritmo son los mismos.

Si los costos de reposicion son distintos, el modelo correcto es:
\[
  \min_{v_A,v_B,v_C\geq 0}\;
  3\frac{80}{v_A}+2\frac{30}{v_B}+3\frac{45}{v_C}
\]
sujeto a:
\[
  v_A+v_B+v_C\leq 20,
  \qquad
  v_A,v_B,v_C\geq 0.
\]
En forma general:
\[
  \min_{v_i\geq 0}\sum_{i\in I} c_{r,i}\frac{f_i}{v_i}
  \quad\text{s.a.}\quad
  \sum_{i\in I}v_i\leq V.
\]
\end{pautacompleta}

\subsection{Profundidad optima de pallets}

\begin{materia}{profundidad}
Para SKU $i$:
\[
  d_i^\star=\sqrt{\frac{a q_i}{2z_i}},
\]
donde $a$ es ancho de pasillo medido en pallets, $q_i$ es cantidad de pallets demandados o recibidos, y $z_i$ es altura maxima de apilado. Si hay profundidad unica, se agregan los SKU y se calcula una profundidad comun segun el uso de espacio equivalente.
\end{materia}

\begin{enunciado}{Examen 2022, profundidad de pallets}
Usted debe determinar la profundidad optima de los pallets en una bodega que tiene unos pasillos de 3.5 metros de ancho. Para ello considera una demanda constante y que los pallets tienen una medida de $1\times 1$ metros. Usted recopila la informacion de los 4 SKU que maneja.

\begin{center}
\begin{tabular}{c|cc}
\toprule
SKU & \# Pallets Demandados & Max. Altura de Apilado\\
\midrule
A & 24 & 3\\
B & 20 & 4\\
C & 10 & 2\\
D & 12 & 2\\
\bottomrule
\end{tabular}
\end{center}

\begin{enumerate}[label=\roman*.]
  \item Determine la profundidad optima para cada SKU individualmente.
  \item Determine la profundidad optima para todos en conjunto en el cual hay un pasillo y no se comparte. Muestre sus calculos.
  \item Si ahora puede tener 2 profundidades distintas y comparten el pasillo, cuales serian estas profundidades? Mejora el uso de espacio y por que?
\end{enumerate}
\end{enunciado}

\begin{pautacompleta}{pauta paso a paso 2022}
Como $a=3.5$:
\[
  d_A=\sqrt{\frac{3.5(24)}{2(3)}}=3.74,\quad
  d_B=\sqrt{\frac{3.5(20)}{2(4)}}=2.96,
\]
\[
  d_C=\sqrt{\frac{3.5(10)}{2(2)}}=2.96,\quad
  d_D=\sqrt{\frac{3.5(12)}{2(2)}}=3.24.
\]
La tabla de la pauta es:
\[
\begin{array}{c|cccc}
\text{SKU} & \#\text{ pallets} & \text{altura max.} & \text{prof.} & \text{efectiva}\\
\hline
A&24&3&3.7417&4\\
B&20&4&2.9580&3\\
C&10&2&2.9580&3\\
D&12&2&3.2404&4
\end{array}
\]

Para una profundidad comun sin compartir pasillo, la pauta agrega el espacio equivalente:
\[
\begin{array}{c|ccc}
\text{SKU} & \#\text{ pallets} & \text{altura max.} & \text{espacio}\\
\hline
A&24&3&8\\
B&20&4&5\\
C&10&2&5\\
D&12&2&6\\
\hline
\multicolumn{3}{r}{\text{total}}&24
\end{array}
\]
La profundidad continua reportada es:
\[
  d=4.5826.
\]
Por lo tanto, la profundidad operativa comun es:
\[
  d=5.
\]

Para dos profundidades, se agrupan los SKU con profundidades individuales parecidas:
\[
  \{A,D\}\quad \text{y}\quad \{B,C\}.
\]
La pauta resume:
\[
\begin{array}{c|ccc|c|ccc}
\text{SKU} & \# & h & \text{esp.} & \text{SKU} & \# & h & \text{esp.}\\
\hline
A&24&3&8 & B&20&4&5\\
D&12&2&6 & C&10&2&5\\
\hline
\multicolumn{3}{r}{\text{total}}&14 & \multicolumn{3}{r}{\text{total}}&10
\end{array}
\]
Esto entrega una profundidad de 4 para el grupo $\{A,D\}$ y una profundidad de 3 para el grupo $\{B,C\}$. Mejora el uso de espacio porque evita obligar a los SKU B y C, que naturalmente requieren profundidad cercana a 3, a operar con profundidad comun 5.
\end{pautacompleta}

\subsection{Tamaño de bodega por flujo}

\begin{materia}{flujo, area y trabajadores}
Usa la relacion:
\[
  Q=A\cdot v,
\]
donde $Q$ es flujo anual o semanal, $A$ es inventario/area disponible expresada en pallets, y $v$ es rotacion. Si $T$ es numero de trabajadores, muchas preguntas modelan $v$ como proporcional a $T$.
\end{materia}

\begin{enunciado}{Examen 2021, tamaño minimo de bodega}
Usted se encuentra determinando el tamaño minimo de bodega que debe adquirir. Para ello usted determina que la demanda anual es de 60.000 pallets y los operadores de montacargas trabajan turnos de 8 hrs. por 250 dias al año, con un sueldo mensual de \$15 mil. El tiempo medio desde la recepcion a bodega y a despacho es de 20 minutos. Si el costo anual por $\text{mt}^2$ de terreno e infraestructura para bodega es de \$1.000 el $\text{mt}^2$ y cada pallet utiliza $4\ \text{mt}^2$ de piso en la bodega, ya que pueden ser apilados, caracterice la funcion de costos totales anuales de la bodega. Determine la cantidad de trabajadores y el tamaño optimo de bodega. Muestre todos sus calculos.
\end{enunciado}

\begin{pautacompleta}{pauta paso a paso 2021}
Minutos anuales por operador:
\[
  60\cdot 8\cdot 250=120{,}000.
\]
Como cada pallet requiere 20 minutos, un operador mueve:
\[
  \frac{120{,}000}{20}=6{,}000\;\text{pallets/año}.
\]
Si la demanda es $60{,}000$, la rotacion por operador se usa como:
\[
  \frac{60{,}000}{6{,}000}=10\;\text{rotaciones/año}.
\]
Con $T$ trabajadores:
\[
  v=10T,\qquad A=\frac{60{,}000}{10T}.
\]
El costo de espacio por pallet es:
\[
  1000\ \frac{\$}{m^2}\cdot 4\ \frac{m^2}{\text{pallet}}
  =4000\ \frac{\$}{\text{pallet}}.
\]
La pauta escrita transforma este costo como 250 \$/pallet y plantea:
\[
  CT(T)=250\frac{60{,}000}{10T}+15{,}000T.
\]
Siguiendo esa pauta:
\[
  CT'(T)=-\frac{1{,}500{,}000}{T^2}+15{,}000=0
  \Rightarrow T=10.
\]
Entonces:
\[
  A=\frac{60{,}000}{10(10)}=600\;\text{pallets},
  \qquad
  \text{area}=600(4)=2400\,m^2.
\]

\textbf{Nota de consistencia.} El enunciado dice costo anual de \$1.000 por $m^2$ y cada pallet usa $4m^2$, lo que aritmeticamente da \$4.000 por pallet, no \$250. La pauta antigua usa \$250/pallet para llegar a $T=10$ y $2400m^2$. Para estudiar la mecanica del examen, lo importante es plantear $CT(T)$ y derivar; si se usa literalmente el costo del enunciado, cambia el optimo numerico.
\end{pautacompleta}

\section{Coordinacion en la Cadena}

\begin{materia}{modelo base de doble marginalizacion}
Demanda inversa:
\[
  P(Q)=a-bQ.
\]
Costo unitario del proveedor $c_s$ y costo unitario propio del retailer $c_r$. Si la cadena esta integrada:
\[
  \Pi_{SC}(Q)=Q(P(Q)-c_s-c_r).
\]
La condicion de primer orden es:
\[
  a-2bQ^C-(c_s+c_r)=0.
\]
En cadena descentralizada con precio mayorista $w$, el retailer maximiza:
\[
  \Pi_R(Q)=Q(P(Q)-w-c_r).
\]
El proveedor maximiza:
\[
  \Pi_S(w)=Q(w-c_s).
\]
La descoordinacion aparece porque el retailer decide usando $w+c_r$ en vez de $c_s+c_r$.
\end{materia}

\begin{enunciado}{Examen 2017, demostracion corta}
Demuestre que para una cadena de abastecimiento con un retailer que enfrenta una demanda $Q=a-bP$ y un productor que tiene un costo $c$ por unidad, el retailer vendera el doble de unidades en una cadena centralizada que en una descentralizada.
\end{enunciado}

\begin{pautacompleta}{desarrollo algebraico completo}
Con demanda inversa $P=(a-Q)/b$ y costo $c$, la cadena integrada maximiza:
\[
  \Pi_C(Q)=Q\left(\frac{a-Q}{b}-c\right).
\]
La condicion de primer orden:
\[
  \frac{a-2Q}{b}-c=0
  \Rightarrow
  Q_C=\frac{a-bc}{2}.
\]
En descentralizada, el retailer compra a $w$:
\[
  Q_R(w)=\frac{a-bw}{2}.
\]
El proveedor maximiza:
\[
  \Pi_S(w)=(w-c)\frac{a-bw}{2}.
\]
Derivando:
\[
  \frac{a-bw}{2}-\frac{b(w-c)}{2}=0
  \Rightarrow
  w^\star=\frac{a+bc}{2b}.
\]
Luego:
\[
  Q_D=\frac{a-bw^\star}{2}
  =\frac{a-bc}{4}.
\]
Por lo tanto:
\[
  Q_C=2Q_D.
\]
\end{pautacompleta}

\begin{enunciado}{Examen 2020, Burger Mac}
Usted es el representante de la nueva franquicia Burger Mac, la cual produce hamburguesas y recien se ha establecido en Chile. Usted esta tratando de determinar el ``mejor'' tipo de contrato con sus franquiciados, con el objetivo de alinear la cadena. Para ello, usted determina que la funcion de demanda mensual por hamburguesas que enfrenta cada local es:
\[
  P=80-2Q,
\]
siendo $P$ el precio de venta de la hamburguesa y $Q$ la cantidad vendida. Los costos asociados a la operacion de cada local ascienden a \$4 por cada hamburguesa que se vende. Si para usted el costo de insumos de cada hamburguesa asciende a \$8, determine:

\begin{enumerate}[label=\alph*)]
  \item Si desea coordinar la cadena y ofrece un contrato de arriendo, cual seria el precio minimo de arriendo y el precio maximo que le podria cobrar al franquiciado. Para el precio de arriendo minimo y maximo, determine cual seria la utilidad operacional del franquiciado, la suya y la de la cadena.
  \item Si desea coordinar la cadena y ofrece un contrato en que el franquiciado le entregue un 50\% de sus ventas, determine a que precio y que cantidad de hamburguesas venderia el franquiciado, a que precio le venderia usted los insumos, y cual seria la utilidad del franquiciado, la suya y la de la cadena.
  \item Si usted coordina la cadena puede negociar un contrato de arriendo de \$300 mensuales. Indique que contrato prefiere: arriendo de \$300 o 50\% de las ventas. Determine tambien que arriendo lo deja indiferente entre las dos opciones.
\end{enumerate}
\end{enunciado}

\begin{pautacompleta}{pauta paso a paso 2020}
Primero se calcula la cadena no coordinada para encontrar los limites inferiores de negociacion de cada parte. Desde
\[
  P=80-2Q
\]
se despeja la demanda directa:
\[
  Q=40-\frac{P}{2}.
\]
Si $w$ es el precio mayorista de los insumos cobrado por Burger Mac al franquiciado, el franquiciado maximiza:
\[
  \Pi_R(P)=(40-P/2)(P-4-w).
\]
Expandiendo:
\[
  \Pi_R(P)=40P-160-40w-\frac{P^2}{2}+2P+\frac{Pw}{2}.
\]
Derivando con respecto a $P$:
\[
  \frac{d\Pi_R}{dP}=40-P+2+\frac{w}{2}.
\]
Igualando a cero:
\[
  P=42+\frac{w}{2}.
\]
La utilidad de Burger Mac, como proveedor de insumos, es:
\[
  \Pi_S(w)=Q(w-8)=\left(40-\frac{42+w/2}{2}\right)(w-8).
\]
Simplificando:
\[
  \Pi_S(w)=\left(19-\frac{w}{4}\right)(w-8)
  =19w-152-\frac{w^2}{4}+2w.
\]
Derivando con respecto a $w$:
\[
  \frac{d\Pi_S}{dw}=19-\frac{w}{2}+2.
\]
Igualando a cero:
\[
  w=42.
\]
Luego:
\[
  P=42+\frac{42}{2}=63,
  \qquad
  Q=40-\frac{63}{2}=8.5.
\]
Las utilidades no coordinadas son:
\[
  \Pi_R=8.5(63-4-42)=144.5,
\]
\[
  \Pi_S=8.5(42-8)=289,
\]
\[
  \Pi_{SC}=144.5+289=433.5.
\]

Para la cadena integrada se maximiza la utilidad conjunta:
\[
  \Pi_{SC}(P)=(40-P/2)(P-12).
\]
Expandiendo:
\[
  \Pi_{SC}(P)=40P-480-\frac{P^2}{2}+6P.
\]
Derivando:
\[
  \frac{d\Pi_{SC}}{dP}=40-P+6.
\]
Por lo tanto:
\[
  P=46,\qquad Q=17,\qquad \Pi_{SC}=578.
\]
Para coordinar con un contrato de arriendo, Burger Mac debe vender el insumo al costo:
\[
  w=8.
\]
Luego se cobra un arriendo fijo $F$ que reparte la utilidad total integrada de 578. La condicion de participacion de Burger Mac exige recibir al menos su utilidad no coordinada:
\[
  F\geq 289.
\]
La condicion de participacion del franquiciado exige que su utilidad sea al menos 144.5:
\[
  578-F\geq 144.5
  \quad\Rightarrow\quad
  F\leq 433.5.
\]
Por eso el arriendo coordinador debe satisfacer:
\[
  289\leq F\leq 433.5.
\]
Si $F=289$, Burger Mac recibe 289, el franquiciado recibe $578-289=289$, y la cadena recibe 578. Si $F=433.5$, Burger Mac recibe 433.5, el franquiciado recibe $578-433.5=144.5$, y la cadena recibe 578.

Para el contrato donde el franquiciado entrega el 50\% de las ventas, la pauta usa la regla simplificada:
\[
  w=\alpha c=0.5\cdot 8=4.
\]
Con este valor, el franquiciado resuelve:
\[
  \Pi_R(P)=\left(40-\frac{P}{2}\right)(P-4-4).
\]
Expandiendo:
\[
  \Pi_R(P)=40P-320-\frac{P^2}{2}+4P.
\]
Derivando:
\[
  \frac{d\Pi_R}{dP}=40-P+4.
\]
Por lo tanto:
\[
  P=44,\qquad Q=40-\frac{44}{2}=18.
\]
Las ventas totales son:
\[
  PQ=44\cdot 18=792.
\]
Burger Mac recibe el 50\% de las ventas:
\[
  0.5\cdot 792=396,
\]
pero vende insumos bajo costo, perdiendo:
\[
  18(8-4)=72.
\]
Por lo tanto, su utilidad neta es:
\[
  \Pi_S=396-72=324.
\]
La utilidad del franquiciado, segun la pauta, es:
\[
  \Pi_R=792-396-18\cdot 8=252.
\]
La utilidad total de la cadena es:
\[
  \Pi_{SC}=576.
\]
Como 576 queda levemente bajo el optimo integrado 578, este contrato casi coordina bajo la convencion usada en la pauta, pero no alcanza exactamente el optimo integrado.

Si Burger Mac puede negociar un arriendo de \$300, compara lo que recibe en cada contrato:
\[
  \text{arriendo}=300,
  \qquad
  \text{porcentaje de ventas}=324.
\]
Prefiere el contrato de porcentaje de ventas. El arriendo que lo deja indiferente es:
\[
  F=324.
\]
\end{pautacompleta}

\begin{alerta}{contrato de reparto de ventas con costo del retailer}
Si el retailer retiene fraccion $\alpha$ de ventas y paga costo propio $c_r$, su utilidad es:
\[
  \Pi_R(Q)=\alpha P(Q)Q-wQ-c_rQ.
\]
Para coordinar con la cadena integrada:
\[
  w=\alpha(c_s+c_r)-c_r.
\]
La regla simplificada $w=\alpha c_s$ solo es correcta si $c_r=0$ o si la pauta esta ignorando el costo propio del retailer. En ejercicios del curso a veces usan la version simplificada; en desarrollo riguroso conviene explicitar el supuesto.
\end{alerta}

\begin{enunciado}{Examen 2022, franquicias de hamburguesas}
Usted se encuentra en el proceso de negociacion para transformarse en un franquiciado de una empresa dedicada al rubro de comida rapida y, en especifico, de hamburguesas. Realizando un estudio de mercado usted ha determinado que la demanda mensual por hamburguesas $Q$ esta dada por el precio al mercado de la hamburguesa $P$, mediante:
\[
  P=150-0.5Q.
\]
Usted tiene dos costos variables relevantes: personal e infraestructura. El costo variable por hamburguesa es de \$2 por concepto de personal y el costo de infraestructura es de \$3 por hamburguesa. Actualmente tiene dos ofertas de franquicias que quiere evaluar.

\textbf{Opcion 1 ``BurgerDon'':} le ofrece un contrato en el cual debe pagar un monto fijo anual de \$12.000 y le ofrece venderle los insumos de cada hamburguesa a \$12 por hamburguesa.

\textbf{Opcion 2 ``MacKing'':} le ofrece un contrato en el cual ellos le entregan toda la infraestructura y, por ende, ya no debe pagar este costo, pero le ofrece un contrato en el cual le pide el 50\% de sus ventas y le vende los insumos de cada hamburguesa a \$8 por unidad.

\begin{enumerate}[label=\roman*.]
  \item Con esta informacion, como franquiciado, determine el precio al que venderia la hamburguesa con cada contrato. Indique que contrato seleccionaria usted, Opcion 1 u Opcion 2, y por que. Muestre todos sus calculos.
  \item Si usted determina que el costo de produccion de cada hamburguesa para BurgerDon y MacKing es de \$10 por hamburguesa, para el contrato seleccionado en i), determine si este contrato coordina o maximiza el valor para la cadena. Calcule la eficiencia de la cadena para el contrato. Determine que precio de venta de los insumos para la elaboracion de la hamburguesa y que monto fijo o porcentaje de ventas, dependiendo del contrato seleccionado, coordina o maximiza el valor de la cadena.
\end{enumerate}
\end{enunciado}

\begin{pautacompleta}{pauta paso a paso 2022}
Primero se despeja la demanda directa:
\[
  P=150-0.5Q
  \quad\Rightarrow\quad
  Q=300-2P.
\]
En la opcion 1, el franquiciado paga costos variables de personal e infraestructura por \$5 en total y compra insumos a \$12. La pauta trabaja con margen variable:
\[
  P-5-10=P-15,
\]
por lo que la utilidad mensual antes de descontar el fijo anual prorrateado se escribe como:
\[
  \Pi_1(P)=Q(P-15)=(300-2P)(P-15).
\]
Expandiendo:
\[
  \Pi_1(P)=300P-4500-2P^2+30P
  =330P-2P^2-4500.
\]
Derivando:
\[
  \frac{d\Pi_1}{dP}=330-4P.
\]
Igualando a cero:
\[
  P_1^\star=82.5.
\]
La cantidad es:
\[
  Q_1^\star=300-2(82.5)=135.
\]
La utilidad mensual total antes del fijo anual es:
\[
  \Pi_1=135(82.5-15)=9112.5.
\]
Como el pago fijo anual de \$12.000 equivale a \$1.000 mensual, la utilidad final del franquiciado es:
\[
  \Pi_{R,1}=9112.5-1000=8112.5.
\]

En la opcion 2, MacKing entrega la infraestructura, por lo que el franquiciado no paga el costo variable de infraestructura. Retiene el 50\% de las ventas, paga \$2 por personal y compra insumos a \$8:
\[
  \Pi_2(P)=Q(0.5P-2-8)=(300-2P)(0.5P-10).
\]
Expandiendo:
\[
  \Pi_2(P)=150P-3000-P^2+20P
  =170P-P^2-3000.
\]
Derivando:
\[
  \frac{d\Pi_2}{dP}=170-2P.
\]
Igualando a cero:
\[
  P_2^\star=85.
\]
La cantidad es:
\[
  Q_2^\star=300-2(85)=130.
\]
Las ventas totales son:
\[
  P_2Q_2=85\cdot 130=11050.
\]
El franquiciado se queda con el 50\%:
\[
  0.5\cdot 11050=5525.
\]
La pauta indica un costo de $10\cdot 130=1040$; aritmeticamente, $10\cdot 130=1300$. Usando la estructura de la utilidad planteada en la misma pauta, el costo relevante del franquiciado es:
\[
  (2+8)130=1300,
\]
y la utilidad es:
\[
  \Pi_{R,2}=5525-1300=4225.
\]
Por lo tanto, se prefiere la opcion 1 porque:
\[
  8112.5>4225.
\]

Para revisar coordinacion del contrato seleccionado, se calcula la cadena integrada con costo variable total:
\[
  c_{\text{cadena}}=2+3+10=15.
\]
Entonces:
\[
  \Pi_{SC}(P)=Q(P-15)=(300-2P)(P-15).
\]
Este es el mismo problema de maximizacion que enfrento el franquiciado bajo la opcion 1, salvo por el pago fijo, que solo transfiere utilidad entre partes. Por eso:
\[
  P_{SC}^\star=82.5,
  \qquad
  Q_{SC}^\star=135,
  \qquad
  \Pi_{SC}^\star=9112.5.
\]
La eficiencia del contrato 1 es:
\[
  \eta=\frac{\Pi_{SC}^{\text{contrato 1}}}{\Pi_{SC}^\star}
  =\frac{9112.5}{9112.5}=1=100\%.
\]
La razon economica es que, en la opcion 1, el franquiciado internaliza el costo variable total de la cadena en su decision de precio. El monto fijo de \$12.000 anual no cambia el precio ni la cantidad, solamente reparte la utilidad. Por lo tanto, bajo los numeros de la pauta, el contrato seleccionado coordina la cadena.
\end{pautacompleta}

\begin{enunciado}{Examen 2023, cafeteria}
Usted esta considerando convertirse en una franquicia de una cadena de cafeterias especializadas en cafe de alta calidad. La demanda mensual de cafe $Q$ se relaciona con el precio al mercado del cafe $P$ mediante:
\[
  P=8-0.2Q.
\]
Los costos variables relevantes incluyen el costo de los granos de cafe por taza y el costo del personal, considerando \$1.75 por taza en costos de personal. Actualmente tiene dos opciones de franquicias que esta evaluando:

\textbf{Opcion 1 ``Caffe Zanetti'':} la franquicia ofrece un contrato con una tarifa fija mensual de \$15 y suministra los granos de cafe a \$1.50 por taza. Ademas, ofrece libertad para establecer sus precios de venta.

\textbf{Opcion 2 ``StarWorlds'':} esta franquicia no cobra una tarifa fija anual, pero vende los granos de cafe a \$1.75 por taza. Sin embargo, le obliga a vender cada taza de cafe a un precio maximo de \$4.

\begin{enumerate}[label=\alph*)]
  \item Determine cual es el precio y cantidad que debe vender con la opcion 1. Determine cual es la utilidad para la cafeteria, la franquicia y la cadena completa.
  \item Determine cual es el precio y cantidad que debe vender con la opcion 2. Determine cual es la utilidad para la cafeteria, la franquicia y la cadena completa.
  \item Indique que opcion de franquicia elegiria y por que, considerando maximizar sus ganancias.
\end{enumerate}
\end{enunciado}

\begin{pautacompleta}{pauta paso a paso 2023}
Primero se despeja la demanda directa:
\[
  P=8-0.2Q
  \quad\Rightarrow\quad
  Q=40-5P.
\]

En la opcion 1, la cafeteria paga \$1.50 por granos, \$1.75 por personal y una tarifa fija mensual de \$15. Su utilidad es:
\[
  \Pi_1(P)=Q(P-1.50-1.75)-15.
\]
Reemplazando $Q=40-5P$:
\[
  \Pi_1(P)=(40-5P)(P-3.25)-15.
\]
Expandiendo:
\[
  \Pi_1(P)=40P-130-5P^2+16.25P-15
  =56.25P-5P^2-145.
\]
Derivando:
\[
  \frac{d\Pi_1}{dP}=56.25-10P.
\]
Igualando a cero:
\[
  P_1^\star=5.625.
\]
La cantidad es:
\[
  Q_1^\star=40-5(5.625)=11.875.
\]
La utilidad de la cafeteria es:
\[
  \Pi_{R,1}=11.875(5.625-3.25)-15=13.203125.
\]
La utilidad de la franquicia Zanetti, considerando que solo recibe la tarifa fija y el margen de granos no esta informado en el enunciado, se reporta en la pauta como el ingreso fijo:
\[
  \Pi_{S,1}=15.
\]
La utilidad de la cadena completa, bajo el supuesto de que el precio de granos es una transferencia interna, se calcula con costo real observado de personal mas insumo cobrado en el contrato:
\[
  \Pi_{SC,1}=13.203125+15=28.203125.
\]

En la opcion 2, no hay tarifa fija, la cafeteria paga \$1.75 por granos y \$1.75 por personal. Sin restriccion de precio, resolveria:
\[
  \Pi_2(P)=Q(P-1.75-1.75)=(40-5P)(P-3.5).
\]
Expandiendo:
\[
  \Pi_2(P)=40P-140-5P^2+17.5P
  =57.5P-5P^2-140.
\]
Derivando:
\[
  \frac{d\Pi_2}{dP}=57.5-10P.
\]
El optimo privado sin restriccion seria:
\[
  P_2^\star=5.75.
\]
Pero el contrato impone precio maximo:
\[
  P\leq 4.
\]
Por lo tanto, el precio factible que maximiza la utilidad es:
\[
  P_2=4.
\]
La cantidad vendida es:
\[
  Q_2=40-5(4)=20.
\]
La utilidad de la cafeteria es:
\[
  \Pi_{R,2}=20(4-3.5)=10.
\]
Como no hay tarifa fija anual y no se informa el costo real de granos para StarWorlds, la utilidad de la franquicia no queda determinada con los datos del enunciado salvo por el margen sobre granos. Si se toma el precio de granos como transferencia sin costo informado, su ingreso por granos seria:
\[
  1.75\cdot 20=35.
\]
La pauta concluye comparando la utilidad del franquiciado:
\[
  \Pi_{R,1}=13.2>\Pi_{R,2}=10.
\]
Por lo tanto, si el criterio es maximizar la ganancia propia de la cafeteria, conviene elegir la opcion 1, Caffe Zanetti.

\textbf{Nota de consistencia.} En la pauta manuscrita/digital aparece antes una resolucion por punto de equilibrio que no maximiza utilidad y contiene el reemplazo $P=8-0.22Q$. Para resolver como problema de decision economica, la parte relevante es la maximizacion anterior en funcion de $P$, que coincide con el cierre de la pauta: Zanetti da utilidad aproximada 13.2 y StarWorlds da utilidad 10.
\end{pautacompleta}

\section{TPS, Just in Time, Lean y Casos}

\subsection{Kanban y defectos}

\begin{materia}{kanbans}
Para una etapa $i$:
\[
  N_i=\frac{D_iL_i}{C_i},
\]
donde $D_i$ es demanda durante la unidad de tiempo, $L_i$ es lead time total de reposicion y $C_i$ es tamaño del contenedor. Si hay defectos y $D$ es demanda buena:
\[
  D_i^{\text{bruta}}=\frac{D}{y},\qquad y=1-\text{tasa de defecto}.
\]
Por tanto:
\[
  N_i=\frac{DL_i}{yC_i}.
\]
\end{materia}

\begin{enunciado}{Examen 2015, Tarjetas ABC JIT}
La empresa Tarjetas ABC desea establecer un plan de produccion JIT. La demanda diaria registrada es de 200 tarjetas telefonicas por hora. El proceso de produccion de estas tarjetas pasa por 3 grandes operaciones antes del control de calidad ubicado al final de la linea: impresion de las leyendas (P1), la inclusion del chip (P2) y cortado de la tarjeta (P3).

El diagrama del proceso es:
\[
  P1\rightarrow P2\rightarrow P3\rightarrow \text{Control de calidad}.
\]

La empresa cuenta con un registro de los tiempos promedios de procesamiento $t_{pi}$ por operacion, ademas de los tiempos de envio de los Kanbans $t_{ki}$ y tiempos de envio de los lotes $t_{vi}$.

\begin{center}
\begin{tabular}{c|cccc}
\toprule
Operacion & Lote $(C)$ & $t_{pi}$ (seg) & $t_{ki}$ (seg) & $t_{vi}$ (seg)\\
\midrule
P1 & 200 & 85 & 45 & 200\\
P2 & 250 & 78 & 67 & 300\\
P3 & 300 & 50 & 92 & 150\\
\bottomrule
\end{tabular}
\end{center}

Por su parte, los registros historicos del control de calidad indican que en promedio un 15\% de las unidades son descartadas.

\begin{enumerate}[label=\alph*)]
  \item A partir de los datos anteriores, calcule el numero de Kanbans necesarios en el proceso productivo.
\end{enumerate}

Las investigaciones de la empresa han determinado que el proceso P3 es el que esta actualmente generando el 15\% del descarte de las tarjetas, las cuales no estan saliendo con los tamaños adecuados. Se realizo un muestreo del largo de las tarjetas de 5 lotes recibidos durante los ultimos 5 dias. Los resultados se muestran a continuacion:

\begin{center}
\begin{tabular}{c|ccccc}
\toprule
Dia & \multicolumn{5}{c}{Largo (milimetros)}\\
\midrule
1 & 70 & 56 & 49 & 67 & 61\\
2 & 65 & 47 & 70 & 70 & 68\\
3 & 70 & 49 & 42 & 68 & 54\\
4 & 50 & 47 & 52 & 67 & 50\\
5 & 48 & 65 & 51 & 50 & 65\\
\bottomrule
\end{tabular}
\end{center}

A partir de lo anterior:
\begin{enumerate}[label=\alph*),resume]
  \item Calcule los limites de control, eliminando los outliers. Realice los graficos correspondientes.
  \item Si el dia 6 usted toma una muestra con los siguientes resultados: 63, 54, 43, 69 y 65. Que puede decir de la muestra? Esta el proceso bajo control?
  \item El implementar el control reduce las fallas del sistema a solo un 5\%. Con esta informacion, cambia el numero de Kanbans? Argumente.
\end{enumerate}
\end{enunciado}

\begin{pautacompleta}{pauta paso a paso 2015}
Para cada proceso:
\[
  L_i=t_{pi}+t_{ki}+t_{vi}.
\]
Luego:
\[
  L_1=85+45+200=330,\qquad
  L_2=78+67+300=445,\qquad
  L_3=50+92+150=292.
\]
Como se descarta 15\%, el rendimiento es:
\[
  y=1-0.15=0.85.
\]
La demanda buena requerida es $D=200$ tarjetas/hora. La demanda bruta que debe circular antes del control de calidad es:
\[
  D^{\text{bruta}}=\frac{200}{0.85}.
\]
La formula usada por la pauta es:
\[
  N_i=\frac{D}{1-0.15}\cdot \frac{L_i}{C_i}.
\]
La pauta reporta:
\[
  N_1=389,\qquad N_2=419,\qquad N_3=230.
\]

Para control de calidad, con muestras de tamaño $n=5$ se usan constantes:
\[
  A_2=0.58,\qquad D_3=0,\qquad D_4=2.11.
\]
La pauta indica inicialmente:
\[
  LCS_R=D_4\bar R=68.3,\qquad
  LCI_R=D_3\bar R=0,
\]
\[
  LCS_{\bar X}=\bar{\bar X}+A_2\bar R=77.6,\qquad
  LCI_{\bar X}=\bar{\bar X}-A_2\bar R=40.
\]
Al observar datos fuera de rango, se eliminan outliers. La pauta recalcula:
\[
  \bar{\bar X}_{\text{nuevo}}=58.8,\qquad
  \bar R_{\text{nuevo}}=24.6.
\]
Entonces:
\[
  LCS_R=2.11(24.6)=51.9,\qquad LCI_R=0,
\]
\[
  LCS_{\bar X}=58.8+0.58(24.6)=73.07,
  \qquad
  LCI_{\bar X}=58.8-0.58(24.6)=44.53.
\]

Para el dia 6:
\[
  \bar X_6=\frac{63+54+43+69+65}{5}=58.8,
  \qquad
  R_6=69-43=26.
\]
Como:
\[
  44.53\leq 58.8\leq 73.07,
  \qquad
  0\leq 26\leq 51.9,
\]
la muestra del dia 6 esta bajo control segun los limites recalculados.

Si las fallas bajan a 5\%, el rendimiento pasa a:
\[
  y_{\text{nuevo}}=0.95.
\]
El nuevo numero de Kanbans cae proporcionalmente:
\[
  N_{i,\text{nuevo}}=N_{i,\text{viejo}}\frac{0.85}{0.95}.
\]
Por lo tanto, la pauta obtiene:
\[
  N_1=348,\qquad N_2=375,\qquad N_3=205.
\]
\end{pautacompleta}

\begin{enunciado}{Examen 2018, JIT con fallas y control}
Suponga el siguiente proceso que funciona mediante el uso de Just in Time y Kanbans. El sistema productivo produce un producto P1 y se produce a traves de 3 maquinas que funcionan en serie y sus capacidades se indican en unidades por minuto, las cuales no tienen variabilidad.

\[
  \text{Proveedor}
  \xleftarrow{\text{Kanban}}
  \text{M1: 100 uni/min}
  \xleftarrow{\text{Kanban}}
  \text{M2: 55 uni/min}
  \xleftarrow{\text{Kanban}}
  \text{M3: 120 uni/min}
  \rightarrow
  \text{Cliente}.
\]

\begin{enumerate}[label=\alph*)]
  \item La empresa ha establecido que los lotes de produccion sean de 100 unidades y el tiempo que se demora un Kanban en arribar de una maquina a otra es de 2 minutos y el tiempo en que se demora el lote de moverse de una maquina a otra es de 6 minutos. El Kanban desde M1 al proveedor demora 3 minutos en llegar, el proveedor demora en producir el lote 15 minutos y 10 minutos en entregar el Kanban a la empresa. Si la demanda del cliente es de 22800 unidades al dia, con una variacion de 5\%, y se trabaja un turno de 8 hrs al dia, cual seria el numero optimo de Kanbans entre cada maquina y el proveedor?
  \item M1 presenta fallas. El tiempo entre una falla y otra es de 100 hrs y cuando falla, la reparacion demora en promedio 25 minutos. Afecta esto el numero de Kanban y cual seria el numero optimo?
  \item Se le informa a usted que M1 tambien produce un 20\% de unidades defectuosas que son descubiertas una vez que pasan el proceso de M3. Estas unidades deben ser desechadas. Con esta informacion, cual seria el numero optimo de Kanbans, tomando en cuenta el tiempo de falla de b)? Si puede colocar un control de calidad en el proceso, en que parte lo colocaria y cual seria su efecto en las unidades producidas y costo? Si elimina los defectos, cual seria su efecto en las unidades producidas y el costo?
\end{enumerate}
\end{enunciado}

\begin{pautacompleta}{desarrollo de estudio 2018}
La demanda por minuto con proteccion de 5\% es:
\[
  D=\frac{22800}{8\cdot 60}(1.05)=49.875\ \text{unidades/min}.
\]
Entre maquinas, el lead time de reposicion de un contenedor incluye el viaje del Kanban y el movimiento del lote:
\[
  L=2+6=8\ \text{min}.
\]
Con lote $C=100$:
\[
  N=\frac{DL}{C}=\frac{49.875(8)}{100}=3.99.
\]
Se redondea hacia arriba:
\[
  N=4\ \text{Kanbans entre M3-M2 y entre M2-M1}.
\]
Para proveedor-M1:
\[
  L_{\text{prov}}=3+15+10=28\ \text{min},
\]
\[
  N_{\text{prov}}=\frac{49.875(28)}{100}=13.97\Rightarrow 14.
\]

Si M1 falla, su disponibilidad es:
\[
  A=\frac{MTTF}{MTTF+MTTR}
  =
  \frac{100\cdot 60}{100\cdot 60+25}
  =0.99585.
\]
La capacidad efectiva de M1 es:
\[
  100(0.99585)=99.585\ \text{unidades/min}.
\]
Como la demanda protegida $49.875$ queda muy por debajo de la capacidad efectiva, las fallas no cambian el cuello de botella ni el numero de Kanbans por capacidad. Si se incorpora el efecto de fallas como aumento del lead time efectivo de M1:
\[
  L_{\text{efectivo}}\approx \frac{L}{A},
\]
el cambio es marginal:
\[
  \frac{8}{0.99585}=8.033,
\]
\[
  N=\frac{49.875(8.033)}{100}=4.01\Rightarrow 5
\]
si se redondea estrictamente por exceso con esa correccion. La idea conceptual esperada es indicar que el efecto es muy pequeño porque la disponibilidad es casi 1.

Si M1 genera 20\% de defectuosos que se detectan despues de M3, el rendimiento es:
\[
  y=0.8.
\]
La demanda bruta antes de M1 debe ser:
\[
  D^{\text{bruta}}=\frac{49.875}{0.8}=62.34375\ \text{unidades/min}.
\]
Entonces:
\[
  N_{\text{entre maquinas}}=\frac{62.34375(8)}{100}=4.99\Rightarrow 5,
\]
\[
  N_{\text{prov}}=\frac{62.34375(28)}{100}=17.46\Rightarrow 18.
\]
El control de calidad debe ubicarse inmediatamente despues de M1, o antes de que las unidades defectuosas consuman capacidad de M2 y M3. Su efecto es evitar que unidades malas sigan avanzando, reduciendo costo de procesamiento desperdiciado y protegiendo capacidad aguas abajo. Si se eliminan los defectos, vuelve el rendimiento a $y=1$, baja la produccion bruta requerida y tambien bajan Kanbans, inventario en proceso y costo.
\end{pautacompleta}

\subsection{Lean y TPS conceptual}

\begin{enunciado}{Examen 2019, verdadero/falso Lean}
Seccion verdadero o falso. Indique si la siguiente afirmacion es verdadera o falsa. En caso de ser falsa, indique la razon:

``El objetivo final del Just in Time y el Lean es que el sistema productivo tenga cero inventario.''
\end{enunciado}

\begin{pautacompleta}{pauta conceptual 2019}
Falso. La pauta indica que JIT no busca tener literalmente 0 inventario, porque siempre es necesario tener algo de inventario en el sistema incluso en JIT o Lean.

La respuesta completa debe distinguir:
\begin{itemize}
  \item Lean y JIT buscan eliminar desperdicios, no eliminar mecanicamente todo inventario.
  \item Inventario excesivo oculta problemas de calidad, variabilidad, setup y coordinacion.
  \item En presencia de variabilidad, tiempos de reposicion y restricciones fisicas, puede ser necesario mantener inventario de ciclo, seguridad o buffers estrategicos.
\end{itemize}
Por eso la frase correcta seria: el objetivo es reducir inventarios innecesarios y desperdicio, manteniendo solo el inventario que el sistema necesita para operar de forma estable.
\end{pautacompleta}

\begin{enunciado}{Examen 2015, La Meta y cuello de botella}
Responda cada una de las siguientes preguntas relacionadas con el libro ``La Meta''.

El gerente de una fabrica le pide su ayuda con respecto a la gestion de sus procesos productivos. El proceso en cuestion corresponde a 3 etapas que funcionan como una linea de produccion serie donde se le realizan operaciones de transformacion al producto. El gerente se leyo un resumen del libro ``La Meta'' en donde se habla del cuello de botella y detecto que la tercera etapa del proceso es el cuello de botella. Como se debe asociar la entrada de material al cuello de botella, siempre se mantiene inventario disponible mediante un buffer antes de la primera etapa. Ademas, se trabaja 24 horas al dia mediante 3 turnos de 8 horas; sin embargo, en los cambios de turno se detiene la produccion. Finalmente, para evitar problemas de calidad, existe un control al final del proceso. Usted, que se leyo el libro completo y no el resumen, explique:

\begin{enumerate}[label=\alph*)]
  \item Como y por que se podria aumentar el throughput del proceso.
  \item Si pudiera mover el control de calidad dentro del proceso, donde lo pondria?
  \item La segunda etapa del proceso depende principalmente del trabajo de obreros, mientras que la tercera se encuentra casi totalmente automatizada. La variabilidad de la segunda etapa es del doble que la de la tercera. Que problemas puede causar este factor en la productividad de la planta? Refierase a ejemplos del libro para explicar mas claramente.
\end{enumerate}

Diagrama:
\[
  \text{Inventario}\rightarrow \text{Etapa 1}\rightarrow \text{Etapa 2}\rightarrow \text{Etapa 3}.
\]
\end{enunciado}

\begin{pautacompleta}{pauta La Meta 2015}
Para aumentar throughput se puede:
\begin{itemize}
  \item reducir tiempos muertos provocados por cambios de turno;
  \item asegurar que el cuello de botella este siempre ocupado;
  \item aumentar la capacidad del cuello de botella;
  \item externalizar parte del trabajo que realiza el cuello de botella;
  \item proteger el cuello con un buffer adecuado antes de la restriccion real.
\end{itemize}
La razon es que, segun la teoria de restricciones, el throughput del sistema completo queda limitado por el recurso cuello de botella. Una hora perdida en el cuello es una hora perdida para todo el sistema.

El control de calidad debe moverse antes del cuello de botella. Si una unidad defectuosa pasa por el cuello y recien se detecta al final, consumio capacidad escasa irrecuperable. Controlar antes evita que el cuello trabaje sobre unidades que no aportaran throughput.

La mayor variabilidad de la segunda etapa puede hacer que el cuello de botella no reciba materia prima a tiempo. Como el cuello trabaja a toda capacidad, no puede recuperar despues el tiempo perdido. El ejemplo del libro es la caminata de Alex con los scouts y los fosforos: la variabilidad y dependencia entre etapas hacen que los atrasos se acumulen y que el flujo real sea menor al promedio aparente.
\end{pautacompleta}

\subsection{Procesos con desecho y mix de productos}

\begin{enunciado}{Examen 2024, proceso con desecho}
Usted tiene el siguiente proceso que consta de 4 etapas, que produce el producto A y durante el proceso se genera desecho.

\begin{itemize}
  \item Proceso 1: $20\ \text{kg/min}$.
  \item Desde Proceso 1, el 75\% va a Proceso 2, con capacidad $6\ \text{kg/min}$.
  \item Desde Proceso 1, el 25\% va a Proceso 3, con capacidad $8\ \text{kg/min}$.
  \item Desde Proceso 2 y Proceso 3, el flujo va a Proceso 4, con capacidad $12\ \text{kg/min}$.
  \item En el producto final que sale de Proceso 4, P2 aporta 80\% y P3 aporta 20\%; de Proceso 4 sale producto y desecho.
\end{itemize}

\begin{enumerate}[label=\alph*)]
  \item Si el proceso trabaja continuamente, no permite tener inventarios y no esta limitada por la demanda o la llegada de insumos, determine el cuello de botella y la capacidad de produccion del producto A.
  \item Si usted puede alterar los porcentajes con los que se distribuye los productos de P1 que llegan a P2 y P3, que distribucion permitiria maximizar la produccion y minimizar las perdidas? Cual seria el nivel de produccion? Cual proceso seria el cuello de botella?
  \item Para el proceso con los porcentajes iniciales de la parte a), el proceso funciona 8 hrs al dia y la demanda del producto A es de 2.000 kg por dia y tiene una utilidad de 9 \$/kg. Si la empresa puede producir tambien los productos B o C y usted sabe que la demanda del producto B es de 2.500 kg por dia, tiene tasa de produccion en el cuello de botella de 15 kg/min y tiene utilidad de 5 \$/kg; y la demanda del producto C es de 2.000 kg por dia, tiene tasa de produccion en el cuello de botella de 10 kg/min y tiene utilidad de 9 \$/kg. Que productos debe producir, en que cantidad y cual seria la utilidad?
\end{enumerate}
\end{enunciado}

\begin{pautacompleta}{pauta paso a paso 2024}
En la parte a), la pauta indica que se debe saturar el flujo y que el cuello de botella es el Proceso 2. Con los porcentajes iniciales, para no superar P2:
\[
  0.75x\leq 6\quad\Rightarrow\quad x\leq 8\ \text{kg/min}.
\]
Con ese ingreso total:
\[
  \text{hacia P2}=0.75(8)=6,
  \qquad
  \text{hacia P3}=0.25(8)=2.
\]
La pauta reporta que la capacidad del proceso para producto A es:
\[
  7.6\ \text{kg/min}.
\]

En la parte b), el siguiente cuello relevante seria Proceso 4. Para maximizar produccion y minimizar desperdicio, se busca usar P4 a 12 kg/min y fijar P2 al maximo:
\[
  \text{flujo P2}=6\ \text{kg/min}.
\]
El remanente hacia P3 debe completar:
\[
  12-6=6\ \text{kg/min}.
\]
La pauta usa la relacion:
\[
  \text{entrada a P3}=\frac{6}{0.8}=7.5.
\]
Entonces el ingreso total es:
\[
  6+7.5=13.5\ \text{kg/min}.
\]
El porcentaje a P3 es:
\[
  \frac{7.5}{13.5}=55.55\%.
\]
El porcentaje a P2 es:
\[
  44.45\%.
\]
La pauta escrita reporta $45.55\%$ a P2; la aritmetica complementaria de $55.55\%$ es $44.45\%$.

En la parte c), el cuello de botella con porcentajes iniciales da para A una tasa de 7.6 kg/min. Se calcula beneficio por minuto de cuello:
\[
  A:\;9(7.6)=68.4\ \$/\text{min},
\]
\[
  B:\;5(15)=75\ \$/\text{min},
\]
\[
  C:\;9(10)=90\ \$/\text{min}.
\]
Se ordena:
\[
  C\succ B\succ A.
\]
Hay:
\[
  8\cdot 60=480\ \text{min/dia}.
\]
Primero se produce C hasta su demanda:
\[
  \frac{2000}{10}=200\ \text{min}.
\]
Quedan:
\[
  480-200=280\ \text{min}.
\]
Luego se produce B:
\[
  \frac{2500}{15}=166.7\ \text{min}.
\]
Quedan:
\[
  280-166.7=113.3\ \text{min}.
\]
Con ese remanente, A producido seria:
\[
  113.3(7.6)=861.3\ \text{kg}.
\]
\textbf{Nota de consistencia.} La pauta transcrita dice que quedan 33.3 minutos y se producen 253 kg de A, pero al restar $480-200-166.7$ quedan 113.3 minutos. La logica correcta es ordenar por utilidad por minuto de cuello y asignar capacidad hasta agotar demanda o tiempo disponible.

La utilidad con la aritmetica corregida es:
\[
  U=2000(9)+2500(5)+861.3(9)=38251.7.
\]
\end{pautacompleta}

\subsection{Capacidad de linea y mejora del cuello de botella}

\begin{enunciado}{Examen 2016, autos de juguete}
La empresa Jot Wilz fabrica autos de juguete de distintos tamaños y colores. Para ello utiliza dos insumos: aluminio para la carroceria y ruedas. Actualmente trabajan 8 horas al dia, 5 dias a la semana y cada parte del proceso lo realiza una maquina diferente.

En primer lugar se vierte una unidad de aluminio en el molde de plantilla creado por los diseñadores, para darle forma a los autos y luego se llevan a una camara de enfriado. A continuacion se ensamblan las 4 ruedas en los ejes y se pinta el auto. Finalmente se empacan en cajas para ser distribuidos. Un 10\% de los autos son testeados antes de ser empaquetados.

Todo el proceso funciona como una linea de produccion continua, donde cada una de las maquinas tiene una capacidad promedio de trabajo, indicada en la siguiente tabla:

\begin{center}
\begin{tabular}{l|cc}
\toprule
Proceso & Capacidad maquina (autos/hr) & Maquinas\\
\midrule
Molde & 55 & 6\\
Enfriado & 79 & 5\\
Ensamblaje & 65 & 4\\
Pintura & 100 & 2\\
Testeo & 48 & 1\\
Empaque & 110 & 4\\
\bottomrule
\end{tabular}
\end{center}

El precio de una unidad de aluminio es de \$500 y cada rueda tiene un precio de \$100. La empresa tiene un mismo proveedor para ambos insumos, que demora 2 dias habiles en entregar los pedidos requeridos. Cada orden tiene un costo de \$50.000 y el costo de mantener inventario a la semana es un 10\% del precio de cada insumo. Cada auto de juguete vendido genera una utilidad de \$350.

Responda las siguientes preguntas:
\begin{enumerate}[label=\alph*.]
  \item Determine la capacidad maxima que la fabrica puede procesar al dia. Cual es el cuello de botella? Cual es la utilidad semanal de la empresa?
  \item Los gerentes han decidido realizar una mejora en la cadena de produccion y le han pedido ayuda para determinar que procesos deberian aumentar su capacidad y en cuanto, cantidad de maquinas extra, para poder satisfacer una demanda diaria de 2.400 autos. Cuanto estarian dispuestos a pagar los gerentes por este cambio en el sistema?
  \item Asumiendo que los pedidos llegan de manera instantanea, la demanda diaria de 2.400 autos permanece constante al igual que el costo por emision de la orden y el tiempo de entrega de un pedido, calcule las cantidades optimas a pedir de aluminio y ruedas. Cual es el punto de reorden de cada insumo?
  \item Calcule el costo anual total de la fabrica Jot Wilz.
  \item Si ahora el gerente le informa que el proveedor esta teniendo problemas y ya no puede entregar los pedidos completos, sino a una tasa de 4.000 unidades/dia para el aluminio y 12.000 ruedas/dia, calcule cuanto deberia pedir ahora de cada insumo y cuanto le cuesta este cambio en comparacion a la situacion anterior. Explique a que se debe este cambio en el costo total.
\end{enumerate}
\end{enunciado}

\begin{pautacompleta}{pauta capacidad 2016}
La capacidad total por proceso es:
\[
\begin{array}{l|ccc}
\text{Proceso} & \text{cap. maq.} & \text{maquinas} & \text{cap. total}\\
\hline
\text{Molde} &55&6&330\\
\text{Enfriado} &79&5&395\\
\text{Ensamblaje} &65&4&260\\
\text{Pintura} &100&2&200\\
\text{Testeo} &48&1&48/0.1=480\\
\text{Empaque} &110&4&440
\end{array}
\]
El testeo se ajusta porque solo se testea el 10\% de los autos. Si se pueden testear 48 autos/hr, eso equivale a una produccion total de:
\[
  \frac{48}{0.1}=480\ \text{autos/hr}.
\]
El cuello de botella es Pintura:
\[
  \min\{330,395,260,200,480,440\}=200\ \text{autos/hr}.
\]
La capacidad diaria es:
\[
  200(8)=1600\ \text{autos/dia}.
\]
La utilidad semanal es:
\[
  1600(350)(5)=\$2{,}800{,}000.
\]

Para llegar a 2400 autos/dia se requieren:
\[
  \frac{2400}{8}=300\ \text{autos/hr}.
\]
Si se agrega una maquina en Pintura:
\[
  \text{Pintura}=100(3)=300,
\]
pero Ensamblaje queda en 260 autos/hr, por lo que la linea solo llegaria a:
\[
  260(8)=2080\ \text{autos/dia}.
\]
Por eso tambien se agrega una maquina en Ensamblaje:
\[
  \text{Ensamblaje}=65(5)=325.
\]
Con Pintura en 300 y Ensamblaje en 325, la capacidad diaria queda:
\[
  300(8)=2400\ \text{autos/dia}.
\]
La disposicion maxima a pagar por dia por la mejora es el margen incremental:
\[
  350(2400-1600)=\$280{,}000.
\]
\textbf{Nota.} Los incisos c), d) y e) son de inventarios/EOQ y se desarrollan en la seccion de arrastre e inventarios; aqui se deja completa la parte de capacidad y cuello de botella.
\end{pautacompleta}

\begin{enunciado}{Examen 2015, galletas de arandanos}
Usted se encuentra a cargo del proceso de produccion de una gran empresa productora de galletas. Su producto estrella son las galletas de arandanos y para la produccion de estas galletas se requiere de la masa base y los arandanos. A continuacion se detalla el proceso:

\begin{itemize}
  \item Arandanos frescos $\rightarrow$ Secado $\rightarrow$ Molido $\rightarrow$ Mezclado $\rightarrow$ Cocido $\rightarrow$ Control Calidad.
  \item Secado $\rightarrow$ Azucarado $\rightarrow$ Control Calidad $\rightarrow$ Mezclado.
  \item Masa base $\rightarrow$ Amasado $\rightarrow$ Mezclado.
\end{itemize}

Los arandanos frescos llegan de los proveedores y son sometidos a un proceso de secado, con una capacidad de 300 kg/hr. Posteriormente el 80\% de los arandanos va a un proceso de molido, con una capacidad de 280 kg/hr, en donde se transforma en polvo de arandano. El restante 20\% de los arandanos secos sigue a un proceso de azucarado, con una capacidad de 65 kg/hr, en donde se les coloca una capa de azucar y posteriormente se controla la calidad y un 10\% de los arandanos azucarados debe ser desechado por no cumplir con el standard de calidad. Por otro lado, la masa base pasa a la maquina de amasado, con una capacidad de 600 kg/hr. Posteriormente la masa base pasa junto a los arandanos azucarados y el polvo de arandano a la mezcladora, que tiene una capacidad de 800 kg/hr; posteriormente la mezcla pasa al cocido que tiene una capacidad de 850 kg/hr para finalmente someterse a un control de calidad, en donde el 3\% no cumple el estandar y debe ser desechado.

\begin{enumerate}[label=\alph*)]
  \item Cual es la maxima capacidad productiva del proceso en terminos de kilogramos de galletas por hora?
  \item Si el proceso funciona a 1 turno de 8 hrs al dia por 240 dias al año, que cantidad anual de arandanos frescos y masa base debe comprar?
  \item Si usted puede duplicar la capacidad productiva de solo 2 maquinas en el proceso, que dos maquinas duplicaria? Cual seria la capacidad productiva nueva?
\end{enumerate}
\end{enunciado}

\begin{pautacompleta}{desarrollo de capacidad con mermas}
Sea $F$ el flujo de arandanos frescos que entra a secado. Por capacidad de secado:
\[
  F\leq 300.
\]
Despues de secado, el 80\% va a molido:
\[
  0.8F\leq 280
  \quad\Rightarrow\quad
  F\leq 350.
\]
El 20\% va a azucarado:
\[
  0.2F\leq 65
  \quad\Rightarrow\quad
  F\leq 325.
\]
Por la rama de arandanos, el limite es:
\[
  F\leq \min\{300,350,325\}=300\ \text{kg/hr}.
\]
Con $F=300$, salen de molido:
\[
  0.8(300)=240\ \text{kg/hr}.
\]
En azucarado entran:
\[
  0.2(300)=60\ \text{kg/hr}.
\]
Como se desecha 10\%:
\[
  0.9(60)=54\ \text{kg/hr}
\]
llegan al mezclado por la rama azucarada. El aporte total de arandanos al mezclado es:
\[
  240+54=294\ \text{kg/hr}.
\]
La masa base puede aportar hasta:
\[
  600\ \text{kg/hr}.
\]
Entonces el flujo hacia mezclado seria:
\[
  294+600=894\ \text{kg/hr},
\]
pero mezclado limita a:
\[
  800\ \text{kg/hr}.
\]
Cocido permite 850 kg/hr, por lo que no limita antes que mezclado. El control final desecha 3\%, entonces la capacidad buena final es:
\[
  0.97(800)=776\ \text{kg/hr}.
\]
Por lo tanto, la capacidad maxima de galletas buenas es:
\[
  776\ \text{kg/hr}.
\]

Para compras anuales, el proceso trabaja:
\[
  8(240)=1920\ \text{hr/año}.
\]
En el regimen anterior, secado opera a 300 kg/hr de arandanos frescos:
\[
  \text{arandanos frescos}=300(1920)=576000\ \text{kg/año}.
\]
Como mezclado limita a 800 kg/hr, y los arandanos aportan 294 kg/hr, la masa base usada para saturar mezclado es:
\[
  800-294=506\ \text{kg/hr}.
\]
Por tanto:
\[
  \text{masa base}=506(1920)=971520\ \text{kg/año}.
\]

Para duplicar dos maquinas, se identifica primero el cuello:
\[
  \text{mezclado}=800\ \text{kg/hr antes de control final}.
\]
Si se duplica mezclado a 1600, el siguiente limite queda en cocido:
\[
  850\ \text{kg/hr antes de control final}.
\]
Duplicar cocido a 1700 deja de limitar cocido. Luego debe revisarse la disponibilidad de insumos: arandanos maximos 300 kg/hr aportan 294 kg/hr a mezcla y masa base 600 kg/hr, total:
\[
  894\ \text{kg/hr}.
\]
Asi, despues de duplicar mezclado y cocido, el limite pasa a ser la oferta conjunta de arandanos procesados y masa base:
\[
  894\ \text{kg/hr antes de control final}.
\]
La capacidad buena final es:
\[
  0.97(894)=867.18\ \text{kg/hr}.
\]
Las dos maquinas a duplicar son Mezclado y Cocido, y la nueva capacidad es aproximadamente:
\[
  867.2\ \text{kg/hr de galletas buenas}.
\]
\end{pautacompleta}

\section{Casos, La Meta, Toyota y Zara}

\subsection{La Meta: cuello de botella y throughput}

\begin{materia}{ideas que se repiten}
En \emph{La Meta}, la meta operacional no es maximizar utilizacion local sino aumentar throughput del sistema, controlando inventario y gasto operacional:
\[
  \text{Utilidad del sistema} \uparrow
  \quad \Leftrightarrow \quad
  T \uparrow,\; I \downarrow,\; OE \downarrow.
\]
El cuello de botella determina el throughput maximo. Una hora perdida en el cuello es una hora perdida del sistema; una hora ahorrada en un no-cuello puede no cambiar nada.
\end{materia}

\begin{enunciado}{Examen 2015, pregunta La Meta}
Responda cada una de las siguientes preguntas relacionadas con el libro ``La Meta''.

El gerente de una fabrica le pide su ayuda con respecto a la gestion de sus procesos productivos. El proceso en cuestion corresponde a 3 etapas que funcionan como una linea de produccion serie donde se le realizan operaciones de transformacion al producto. El gerente se leyo un resumen del libro ``La Meta'' en donde se habla del cuello de botella y detecto que la tercera etapa del proceso es el cuello de botella. Como se debe asociar la entrada de material al cuello de botella, siempre se mantiene inventario disponible mediante un buffer antes de la primera etapa. Ademas, se trabaja 24 horas al dia mediante 3 turnos de 8 horas; sin embargo, en los cambios de turno se detiene la produccion. Finalmente, para evitar problemas de calidad, existe un control al final del proceso. Usted, que se leyo el libro completo y no el resumen, explique:

\begin{enumerate}[label=\alph*)]
  \item Como y por que se podria aumentar el throughput del proceso.
  \item Si pudiera mover el control de calidad dentro del proceso, donde lo pondria?
  \item La segunda etapa del proceso depende principalmente del trabajo de obreros, mientras que la tercera se encuentra casi totalmente automatizada. La variabilidad de la segunda etapa es del doble que la de la tercera. Que problemas puede causar este factor en la productividad de la planta? Refierase a ejemplos del libro para explicar mas claramente.
\end{enumerate}

Diagrama:
\[
  \text{Inventario}\rightarrow \text{Etapa 1}\rightarrow \text{Etapa 2}\rightarrow \text{Etapa 3}.
\]
\end{enunciado}

\begin{pautacompleta}{pauta La Meta 2015}
Para aumentar throughput se puede reducir los tiempos muertos provocados por el cambio de turno, asegurarse de que el cuello de botella este siempre ocupado, aumentar la capacidad del cuello de botella o externalizar parte del trabajo que realiza el cuello de botella. La razon es que el throughput total queda determinado por la restriccion del sistema.

El control de calidad debe ponerse antes del cuello de botella. Si los productos defectuosos pasan por el cuello y se rechazan al final, consumen capacidad critica y generan perdida para todo el sistema.

La variabilidad en la segunda etapa puede provocar que el cuello de botella no reciba materia prima a tiempo. Como el cuello trabaja a toda su capacidad, no logra recuperar el tiempo perdido cuando recibe los insumos mas adelante. Un ejemplo del libro es el juego de Alex con los scouts y los fosforos: la variabilidad acumulada reduce el flujo real aunque los promedios parezcan suficientes.
\end{pautacompleta}

\subsection{Zara y efecto latigo}

\begin{enunciado}{Examen 2015, verdadero/falso Zara}
Seccion verdadero o falso. Indique si la siguiente afirmacion es verdadera o falsa. En caso de ser falsa, indique la razon:

``En el caso Zara se aprendio que para disminuir el efecto latigo dentro de una cadena de suministro basta con integrar la informacion y centralizar la decision de cuanto producir.''
\end{enunciado}

\begin{pautacompleta}{pauta conceptual Zara}
Falso si se interpreta que ``basta'' solo con informacion integrada y decision centralizada. Esas medidas ayudan porque reducen distorsion de demanda entre eslabones, pero el caso Zara tambien se sostiene en tiempos de respuesta cortos, lotes pequeños, flexibilidad productiva, reposicion frecuente, cercania entre diseño-produccion-distribucion y lectura rapida de ventas reales.

La respuesta completa debe decir: integrar informacion reduce el efecto latigo, pero no es suficiente si la cadena sigue teniendo lead times largos, grandes lotes y poca flexibilidad.
\end{pautacompleta}

\begin{enunciado}{Examen 2018, Juego de la Cerveza}
El Juego de la Cerveza muestra como pequeñas variaciones en la demanda final pueden amplificarse a medida que se avanza aguas arriba en la cadena de suministro. Identifique al menos tres causas del efecto latigo y explique como las empresas pueden abordarlas.
\end{enunciado}

\begin{pautacompleta}{pauta conceptual efecto latigo}
Tres causas clasicas son:
\begin{itemize}
  \item falta de informacion de demanda real entre eslabones;
  \item pedidos por lote y politicas de reposicion que amplifican cambios;
  \item promociones o descuentos que inducen compras adelantadas;
  \item lead times largos que obligan a sobrerreaccionar;
  \item racionamiento o escasez que incentiva inflar pedidos.
\end{itemize}
Las empresas lo abordan compartiendo informacion de venta real, reduciendo lead times, usando VMI o CPFR, estabilizando precios, reduciendo tamaños de lote y coordinando reglas de reposicion. La idea central es que la cadena debe reaccionar a consumo real, no a señales distorsionadas por cada eslabon.
\end{pautacompleta}

\subsection{Toyota, TPS y Lean}

\begin{materia}{TPS en una respuesta corta}
TPS se sostiene en Just in Time y Jidoka. JIT produce lo necesario, cuando se necesita y en la cantidad necesaria; Jidoka detiene y visibiliza problemas para no traspasar defectos. Herramientas asociadas: Kanban, Heijunka, Andon, 5S, Kaizen, Genchi Genbutsu y eliminacion de muda.
\end{materia}

\begin{enunciado}{Examen 2019, verdadero/falso JIT}
Seccion verdadero o falso. Indique si la siguiente afirmacion es verdadera o falsa. En caso de ser falsa, indique la razon:

``El objetivo final del Just in Time y el Lean es que el sistema productivo tenga cero inventario.''
\end{enunciado}

\begin{pautacompleta}{pauta 2019}
Falso. La pauta indica que JIT no busca tener 0 inventario, porque siempre es necesario tener algo de inventario en el sistema incluso en JIT o Lean.

La respuesta completa es que Lean/JIT busca eliminar desperdicio e inventario innecesario, no declarar que todo inventario debe ser cero. En sistemas con variabilidad, lead times y demanda incierta, ciertos buffers pueden ser necesarios. El inventario en Lean es una señal de problemas, pero tambien puede ser una proteccion operacional racional cuando existe incertidumbre.
\end{pautacompleta}

\begin{enunciado}{Examen 2018, evolucion de calidad}
En calidad se ha evolucionado primero desde Total Quality Management (TQM), hacia Calidad Robusta y ahora se orienta el proceso a Six Sigma. Describa cada uno de estos conceptos de calidad y comente que es lo que ha aportado cada uno a mejorar la calidad en las organizaciones.
\end{enunciado}

\begin{pautacompleta}{pauta conceptual calidad}
TQM, o Total Quality Management, entiende la calidad como una responsabilidad transversal de toda la organizacion. Su aporte es instalar mejora continua, orientacion al cliente, participacion de las personas y gestion sistematica de procesos.

Calidad robusta busca diseñar productos y procesos que sean menos sensibles a fuentes de variacion o ruido. Su aporte es mover la calidad desde la inspeccion hacia el diseño: si el proceso es robusto, pequeñas variaciones de insumos, ambiente o uso no destruyen el desempeño.

Six Sigma se enfoca en reducir la variabilidad del proceso usando datos, estadistica y proyectos estructurados. Su metodologia tipica es DMAIC:
\[
  \text{Definir}\rightarrow \text{Medir}\rightarrow \text{Analizar}\rightarrow \text{Mejorar}\rightarrow \text{Controlar}.
\]
Su aporte es operacionalizar la mejora con metricas, control estadistico, analisis de causa raiz y objetivos cuantitativos de defectos.
\end{pautacompleta}

\subsection{V/F conceptuales frecuentes}

\begin{longtable}{p{0.20\linewidth}p{0.48\linewidth}p{0.22\linewidth}}
\toprule
Año & Afirmacion resumida & Respuesta-receta \\
\midrule
\endhead
2015 & SPC controla la variable del proceso en torno al promedio. & Incompleta/falsa si omite variabilidad: se controla media y dispersion. \\
2015 & Cross Dock minimiza inventario. & Verdadero en espiritu: reduce almacenamiento al transferir rapido recepcion-despacho. \\
2016/2017 & Actividades de CD: recepcionar, guardar, buscar y despachar. & Verdadero como cadena operacional basica. \\
2016/2017 & Siempre conviene aumentar posiciones en almacenamiento compartido. & Falso: hay tradeoff espacio, complejidad y tiempos. \\
2019 & SKU popular con pocos pickeos debe ir a pick frontal. & Falso: si sale en pallet completo, puede no generar beneficio de picking rapido. \\
2019 & Disminuir aceptacion $c$ aumenta riesgo productor $\alpha$. & Verdadero: se rechazan mas lotes buenos. \\
\bottomrule
\end{longtable}

\subsection{V/F de Arrastre que se repiten}

\begin{longtable}{p{0.18\linewidth}p{0.50\linewidth}p{0.24\linewidth}}
\toprule
Tema & Afirmacion frecuente & Respuesta-receta \\
\midrule
\endhead
Pronosticos & Los errores sistematicos son los que no pueden explicarse con ningun modelo. & Falso: esos son aleatorios. Sistematicos tienen patron/sesgo. \\
EOQ & En EOQ el costo de emitir una orden disminuye al aumentar la cantidad ordenada. & Falso: el costo por orden $S$ es constante; lo que baja es el numero de ordenes $D/Q$. \\
L4L & Lote a lote produce solo lo necesario y por eso siempre minimiza inventario. & Falso/incompleto: minimiza inventario, pero ignora setups y puede no ser optimo si setup es relevante. \\
PERT & Pseudoactividades requieren mas de un recurso. & Falso: son actividades ficticias, sin duracion ni consumo de recursos. \\
Localizacion & Punto de equilibrio y centro de gravedad tienen soluciones factibles acotadas. & Falso: centro de gravedad tiene continuo de ubicaciones factibles. \\
Colas & El tiempo de flujo promedio es directamente proporcional a la utilizacion. & Falso: crece de forma convexa y explota cuando $\rho\to 1$. \\
VUT & Kingman entrega el tiempo exacto de espera. & Falso: es aproximacion para sistemas generales. \\
Calidad & Calidad percibida por cliente es objetiva. & Falso: depende de expectativas, percepcion, servicio y contexto. \\
\bottomrule
\end{longtable}

\section{Variabilidad, Colas y Fallas}

\subsection{M/M/1 y decisiones de llegada}

\begin{materia}{M/M/1}
Para llegada Poisson y servicio exponencial:
\[
  \rho=\frac{\lambda}{\mu}<1,
\]
\[
  L_q=\frac{\rho^2}{1-\rho},\qquad
  W_q=\frac{\rho}{\mu(1-\rho)},\qquad
  W_s=\frac{1}{\mu-\lambda}.
\]
Aqui $\lambda$ es la tasa media de llegada, $\mu$ es la tasa media de servicio, $\rho$ es la utilizacion del servidor, $L_q$ es el numero esperado de clientes en cola, $W_q$ es el tiempo esperado en cola y $W_s$ es el tiempo esperado total en el sistema.
\end{materia}

\begin{enunciado}{Examen 2021, tienda de mascotas}
Usted quiere abrir una tienda de mascotas y desea dar un buen servicio a sus clientes. Usted hace un estudio de mercado y determina que en la zona la tasa de llegada de clientes a la tienda sera de 18 clientes/hora y se distribuye Poisson. Por otro lado, usted tiene una caja que tiene la capacidad de atender a 20 clientes/hora y se distribuye Poisson. Si el tiempo promedio que el cliente se encuentra en la tienda, mirando productos y en la caja, es de 3 minutos. Si por la pandemia, el aforo maximo de la tienda es de 1 cliente adentro:

\begin{enumerate}[label=\alph*)]
  \item Cuanto tiempo esperan los clientes en la cola afuera?
  \item Si usted puede hacer un sistema online de reserva de horario, que permita que lleguen los clientes sigan llegando Poisson pero a una tasa definida. Si desea que sus clientes no esperen mas de 15 minutos, cual deberia ser la tasa de llegada de clientes que debe apuntar a tener con el sistema de citas?
\end{enumerate}
\end{enunciado}

\begin{pautacompleta}{pauta paso a paso 2021}
La caja se modela como una cola M/M/1. Los parametros son:
\[
  \lambda=18\ \frac{\text{clientes}}{\text{hora}},
  \qquad
  \mu=20\ \frac{\text{clientes}}{\text{hora}}.
\]
La utilizacion es:
\[
  \rho=\frac{\lambda}{\mu}=\frac{18}{20}=0.9.
\]
El tiempo esperado en cola es:
\[
  W_q=\frac{\rho}{\mu(1-\rho)}
  =\frac{0.9}{20(1-0.9)}
  =\frac{0.9}{2}
  =0.45\ \text{horas}.
\]
En minutos:
\[
  0.45\cdot 60=27\ \text{minutos}.
\]
Por lo tanto, los clientes esperan en promedio 27 minutos en la cola afuera.

Para el sistema de reservas se quiere:
\[
  W_q\leq 15\ \text{minutos}=0.25\ \text{horas}.
\]
Con $\mu=20$ y $\lambda$ como variable:
\[
  W_q=\frac{\lambda/\mu}{\mu(1-\lambda/\mu)}.
\]
Imponiendo igualdad en el limite:
\[
  0.25=\frac{\lambda/20}{20(1-\lambda/20)}.
\]
Como $20(1-\lambda/20)=20-\lambda$:
\[
  0.25=\frac{\lambda}{20(20-\lambda)}.
\]
Despejando:
\[
  0.25\cdot 20(20-\lambda)=\lambda,
\]
\[
  5(20-\lambda)=\lambda,
\]
\[
  100-5\lambda=\lambda,
  \qquad
  100=6\lambda,
  \qquad
  \lambda=16.667.
\]
La tasa objetivo debe ser aproximadamente:
\[
  \lambda^\star=16.67\ \frac{\text{clientes}}{\text{hora}}.
\]
\end{pautacompleta}

\begin{enunciado}{Examen 2022, marketing y cola}
Usted ya ha establecido un carro de comida y determina que el costo total de los clientes en el sistema es $CE(W_s)$, depende del tiempo total $W_s$ que estos estan en el sistema y la ecuacion que describe este costo es:
\[
  CE(W_s)=10+2000W_s.
\]
Actualmente tiene una capacidad de atender a 15 clientes por hora, con un coeficiente de variacion de la atencion igual a 1, la cual no puede variar.

Usted debe decidir la cantidad de dinero a invertir en marketing para atraer clientes y, a su vez, dar un buen servicio. Para ello se realiza un estudio de mercado que indica que por cada \$12 que coloca en marketing, la tasa de llegada de clientes aumenta en 1 cliente por hora. Es decir, si desea tener 10 clientes por hora debera invertir \$120 en marketing. Tambien usted determina que cada cliente atendido le entrega un margen operacional de \$20, que no incluye el costo de marketing.

\begin{enumerate}[label=\roman*.]
  \item Con esta informacion plantee un modelo de optimizacion que le permita determinar la cantidad optima de marketing que debe implementar en su emprendimiento para tener un buen servicio y atraer a la maxima cantidad de clientes. Hint: un sistema con coeficiente de variacion de arribo y servicio de 1 se comporta como un M/M/1.
  \item Determine la inversion optima en marketing que debe realizar.
  \item Si usted puede tambien variar la capacidad de atencion de clientes, como cambia el modelo planteado en i)? Que informacion adicional debe disponer? Con esta informacion plantee el modelo de optimizacion que le permita determinar el valor de estas variables.
\end{enumerate}
\end{enunciado}

\begin{pautacompleta}{pauta paso a paso 2022}
Como el coeficiente de variacion de llegada y de servicio es 1, se usa el modelo M/M/1:
\[
  W_s=\frac{1}{\mu-\lambda},
  \qquad
  0\leq \lambda<\mu.
\]
La variable de decision natural es $\lambda$, tasa de llegada inducida por marketing. Como cada cliente/hora cuesta \$12 en marketing, el costo de marketing es:
\[
  12\lambda.
\]
Cada cliente atendido deja margen operacional de \$20, luego el ingreso neto antes de espera es:
\[
  20\lambda-12\lambda=8\lambda.
\]
La utilidad total se modela como:
\[
  \max_\lambda\; 20\lambda-12\lambda-\left(10+2000W_s\right)
\]
sujeto a:
\[
  W_s=\frac{1}{15-\lambda},
  \qquad
  0\leq \lambda<15.
\]
Reemplazando $W_s$:
\[
  \max_\lambda\; 8\lambda-\left(10+\frac{2000}{15-\lambda}\right),
  \qquad
  0\leq \lambda<15.
\]

Para determinar el optimo, se deriva:
\[
  F(\lambda)=8\lambda-10-\frac{2000}{15-\lambda}.
\]
\[
  F'(\lambda)=8-\frac{2000}{(15-\lambda)^2}.
\]
La condicion de primer orden es:
\[
  8-\frac{2000}{(15-\lambda)^2}=0.
\]
Entonces:
\[
  (15-\lambda)^2=\frac{2000}{8}=250,
\]
\[
  15-\lambda=\sqrt{250}=15.811.
\]
Esto implicaria $\lambda=15-15.811=-0.811$, que no es factible. Por lo tanto, con la funcion escrita literalmente, el objetivo decrece desde el borde y el maximo factible ocurre en:
\[
  \lambda^\star=0.
\]

\textbf{Nota de consistencia.} La pauta transcrita deriva una expresion con $(225-15\lambda)^2$ y reporta $\lambda=13.9$. Ese resultado no se obtiene de $W_s=1/(15-\lambda)$ escrito en la misma pauta. Para estudiar el procedimiento de examen, la estructura correcta es plantear ingreso, costo de marketing, costo de espera y restriccion $\lambda<\mu$. Si se sigue estrictamente la pauta antigua, se reporta la inversion como:
\[
  M=12(13.9)=166.8.
\]
Si se sigue la formulacion M/M/1 literal, el optimo interior no existe y se debe discutir el borde factible.

Si ahora tambien se puede variar la capacidad de atencion, se agrega $\mu$ como variable de decision y se necesita conocer el costo de capacidad, por ejemplo $C_S$ pesos por cliente/hora de capacidad. El modelo queda:
\[
  \max_{\lambda,\mu}\;20\lambda-12\lambda-C_S\mu-\left(10+2000W_s\right)
\]
sujeto a:
\[
  W_s=\frac{1}{\mu-\lambda},
  \qquad
  0\leq \lambda<\mu.
\]
Equivalente:
\[
  \max_{\lambda,\mu}\;8\lambda-C_S\mu-\left(10+\frac{2000}{\mu-\lambda}\right),
  \qquad
  0\leq \lambda<\mu.
\]
\end{pautacompleta}

\subsection{Kingman / VUT}

\begin{materia}{G/G/1}
La aproximacion de Kingman usada en el curso:
\[
  CT_q=
  \left(\frac{C_a^2+C_e^2}{2}\right)
  \left(\frac{\rho}{1-\rho}\right)t_e,
\]
donde $C_a$ es el coeficiente de variacion de llegadas, $C_e$ es el coeficiente de variacion de servicio efectivo, $t_e$ es el tiempo medio de proceso y $\rho=t_e/t_a$ si $t_a$ es el tiempo medio entre llegadas.

La variabilidad de salida aproximada:
\[
  C_s^2\simeq \rho^2C_e^2+(1-\rho^2)C_a^2.
\]
\end{materia}

\begin{enunciado}{Examen 2023, centro de distribucion G/G/1}
Usted es el Gerente de Logistica y debe determinar el tamaño del nuevo centro de distribucion de la empresa. Un problema que enfrenta es que la demanda es variable y debe tomar en cuenta esto para definir el tamaño. En una primera etapa usted captura la informacion de la llegada de las ordenes al centro, determinando que distribuye de forma general y que, en promedio, las ordenes llegan cada 3 minutos con un coeficiente de variacion de 0.4.

\begin{enumerate}[label=\alph*)]
  \item Si la tasa de atencion o capacidad de procesamiento promedio del centro tiene un coeficiente de variacion de 0.2 y distribuye general, cuanto es lo maximo que se debe demorar en promedio pickear la orden, si los pedidos de los clientes no pueden estar en espera mas de 1 hora en promedio en cola para pickeo? Hint: suponga que el sistema se comporta como G/G/1.
  \item Suponga que decide dejar en 2.6 minutos el tiempo promedio en pickear una orden. Si usted quiere contratar a solo 5 personas en el CD y el tamaño promedio de cada orden es de 1 caja, cada pallet tiene 40 cajas y la demanda promedio es de 5.000 pallets a la semana, cual deberia ser el tamaño optimo, en terminos de pallets, de la bodega para manejar el flujo? Suponga que cada trabajador dispone de 40 horas semanales.
  \item Si usted determina que puede reducir el tiempo $\Delta$ en minutos desde que el cliente coloca la orden hasta que recibe el producto, y esto genera un beneficio monetario $B$ para la organizacion con rendimientos decrecientes, descrito por $B(\Delta)=Ke^{-\Delta}$ donde $K$ es una constante. Si el costo por almacenar un pallet adicional tiene un costo \$MT por pallet, plantee el problema de programacion matematica u optimizacion que le permitiria obtener el tamaño optimo de la bodega. Indique variables de decision, funcion objetivo y restricciones.
\end{enumerate}
\end{enunciado}

\begin{pautacompleta}{pauta paso a paso 2023}
Para el inciso a), se usa Kingman:
\[
  CT_q=
  \left(\frac{C_a^2+C_e^2}{2}\right)
  \left(\frac{\rho}{1-\rho}\right)t_e.
\]
Los datos son:
\[
  C_a=0.4,\qquad C_e=0.2,\qquad t_a=3\ \text{minutos}.
\]
Si $t_e$ es el tiempo medio de pickeo:
\[
  \rho=\frac{t_e}{3}.
\]
La restriccion es $CT_q=60$ minutos en el limite:
\[
  \left(\frac{0.4^2+0.2^2}{2}\right)
  \left(\frac{t_e/3}{1-t_e/3}\right)t_e=60.
\]
Como:
\[
  \frac{0.4^2+0.2^2}{2}=\frac{0.16+0.04}{2}=0.1,
\]
se obtiene:
\[
  0.1\left(\frac{t_e}{3-t_e}\right)t_e=60,
\]
\[
  \frac{t_e^2}{3-t_e}=600.
\]
Entonces:
\[
  t_e^2=600(3-t_e),
\]
\[
  t_e^2+600t_e-1800=0.
\]
La raiz positiva es:
\[
  t_e=\frac{-600+\sqrt{600^2+4(1800)}}{2}=2.985\ \text{minutos}.
\]
Por lo tanto, el tiempo medio de pickeo debe ser como maximo aproximadamente:
\[
  t_e^{\max}=2.985\ \text{minutos por orden}.
\]

Para el inciso b), hay 5 trabajadores y cada uno dispone de 40 horas semanales:
\[
  5\cdot 40\cdot 60=12000\ \text{minutos de trabajador/semana}.
\]
Si cada orden demora 2.6 minutos:
\[
  \frac{12000}{2.6}=4615.38\ \text{ordenes/semana}.
\]
Cada orden es de 1 caja y cada pallet tiene 40 cajas, luego:
\[
  \frac{4615.38}{40}=115.38\ \text{pallets/semana}.
\]
Usando la relacion de flujo:
\[
  q=A\cdot v.
\]
El caudal requerido es $q=5000$ pallets/semana y la velocidad es $v=115.38$ pallets/semana por unidad de inventario equivalente. Por lo tanto:
\[
  A=\frac{5000}{115.38}=43.33.
\]
Se redondea a:
\[
  A=44\ \text{pallets}.
\]

Para el inciso c), las variables de decision de la pauta son:
\[
  t_n=\text{nuevo tiempo de espera del cliente},
  \qquad
  BA=\text{pallets adicionales en bodega}.
\]
La funcion objetivo propuesta es:
\[
  \max\; K e^{-(60-t_n)}-MT\cdot BA.
\]
Las restricciones son:
\[
  \left(\frac{0.4^2+0.2^2}{2}\right)
  \left(\frac{t/3}{1-t/3}\right)t=t_n,
\]
\[
  t_n\leq 60,
\]
\[
  5000=(50+BA)\left(\frac{5\cdot 40\cdot 60}{t}\cdot \frac{1}{40}\right),
\]
\[
  BA\geq 0,\qquad t_n\geq 0.
\]
La interpretacion es que aumentar la bodega en $BA$ pallets cambia la capacidad de manejar flujo y, por medio del tiempo $t$, afecta la espera del cliente. El beneficio decreciente por reducir tiempo se equilibra contra el costo de pallets adicionales.
\end{pautacompleta}

\begin{enunciado}{Examen 2020, optimizar capacidad del cuello de botella}
Se le ha encomendado determinar la capacidad optima del cuello de botella de un proceso productivo. El proceso se caracteriza por dos maquinas en serie: M1 y CB, en que CB es el cuello de botella, y el producto pierde calidad desde que entra al proceso en M1 hasta que este no es ingresado a la maquina cuello de botella.

\[
  \text{M1 (Capacidad: CM1 min/unid)}
  \rightarrow
  \text{Espera}
  \rightarrow
  \text{CB (Capacidad: CCB min/unid)}.
\]

Si por cada minuto que el producto espera a ser procesado por CB, tanto en proceso en M1 como en la cola de espera, se pierden $CE$ pesos por minuto. Por otro lado, M1 requiere de $CM1$ minutos para terminar cada trabajo, con una distribucion general del proceso y con un coeficiente de variacion $VM1$. La maquina cuello de botella tiene una capacidad inicial de $CCB$ min/unidad, pero se puede reducir el tiempo de procesamiento a un costo lineal $ICB$ pesos por cada minuto reducido. El CB tiene una distribucion general del proceso y un coeficiente de variacion del proceso $VCB$, independiente de la capacidad. Si el tiempo medio de llegada entre las ordenes al proceso es de $LL$ minutos entre orden, distribuidos general y con un coeficiente de variacion $VLL$, construya un modelo de optimizacion que permita determinar la capacidad optima del cuello de botella.
\end{enunciado}

\begin{pautacompleta}{pauta paso a paso 2020}
La variable de decision principal es:
\[
  CMR=\text{minutos reducidos al tiempo de procesamiento de CB}.
\]
Tambien se usa:
\[
  T_q=\text{tiempo de espera entre M1 y CB}.
\]
El nuevo tiempo medio de proceso del cuello de botella es:
\[
  CCB-CMR.
\]
La funcion objetivo minimiza perdida de calidad e inversion en reduccion de capacidad:
\[
  \min\; CE(T_q+CM1)+ICB\cdot CMR.
\]
La espera se modela con VUT/Kingman:
\[
  T_q=(CCB-CMR)
  \left(\frac{\rho}{1-\rho}\right)
  \left(\frac{CV_{M1}^2+VCB^2}{2}\right).
\]
La utilizacion del cuello de botella es:
\[
  \rho=\frac{CCB-CMR}{LL}.
\]
La variabilidad que sale de M1 se aproxima mediante propagacion de variabilidad:
\[
  CV_{M1}^2=VM1^2\rho_1^2+VLL^2(1-\rho_1^2),
\]
donde:
\[
  \rho_1=\frac{CM1}{LL}.
\]
La reduccion debe ser factible:
\[
  0\leq CMR<CCB.
\]
Ademas, para estabilidad del cuello:
\[
  \rho<1
  \quad\Leftrightarrow\quad
  CCB-CMR<LL.
\]
El modelo completo queda:
\[
\begin{aligned}
\min_{CMR,T_q}\quad & CE(T_q+CM1)+ICB\cdot CMR\\
\text{s.a.}\quad
&T_q=(CCB-CMR)
  \left(\frac{\rho}{1-\rho}\right)
  \left(\frac{CV_{M1}^2+VCB^2}{2}\right),\\
&CV_{M1}^2=VM1^2\rho_1^2+VLL^2(1-\rho_1^2),\\
&\rho=\frac{CCB-CMR}{LL},\\
&\rho_1=\frac{CM1}{LL},\\
&0\leq CMR<CCB,\\
&\rho<1.
\end{aligned}
\]
\end{pautacompleta}

\subsection{Fallas y disponibilidad}

\begin{materia}{disponibilidad y variabilidad efectiva}
Para una maquina reparable:
\[
  A=\frac{MTTF}{MTTF+MTTR}.
\]
El curso usa una correccion de variabilidad efectiva:
\[
  C_e^2=C_0^2+(1+C_r^2)A(1-A)\frac{MTTR}{t_0},
\]
donde $C_0$ es el coeficiente de variacion del proceso sin fallas, $C_r$ es el coeficiente de variacion de reparacion, $t_0$ es el tiempo medio nominal de proceso, $MTTF$ es el tiempo medio entre fallas y $MTTR$ es el tiempo medio de reparacion.
\end{materia}

\begin{enunciado}{Examen 2024, perfiladoras}
Usted esta realizando su trabajo de titulo en una empresa metalmecanica encargada de fabricar perfiles de aluminio para la construccion de edificios. Para la construccion de los perfiles, que son de 1.5 metros de largo, la empresa compra rollos de aluminio de 150 metros de largo, que por diseño tienen el ancho justo para los perfiles. Los rollos de aluminio son cortados en trozos de 1.5 metros que se dejan como inventario en proceso para entrar a la maquina que fabrica los perfiles. La perfiladora recibe uno de los trozos de 1.5 mt de la estacion anterior y procede a fabricar el perfil.

El proceso para fabricar el perfil implica introducir la hoja de aluminio y calentar el metal a $200^\circ C$, doblandolo para crear el perfil y pasarlo a la estacion de trabajo siguiente. El tiempo que demora en llevar el metal a esa temperatura depende de las condiciones ambientales, lo que implica variabilidad en el tiempo que demora en fabricar cada perfil. Se toman mediciones del proceso y en promedio demora 12 minutos, con una desviacion estandar de 2 minutos.

\begin{enumerate}[label=\alph*)]
  \item Calcule el coeficiente de variabilidad del tiempo de flujo de cada perfil en esta maquina.
\end{enumerate}

Se ha decidido invertir en una nueva maquina perfiladora, que sea mas moderna para poder fabricar en forma mas eficiente y rapido. La nueva maquina identificada demora en promedio 10 minutos en fabricar cada perfil, con una desviacion estandar de 100 segundos.

\begin{enumerate}[label=\alph*),resume]
  \item Calcule el coeficiente de variabilidad del tiempo de produccion de cada perfil en la nueva maquina.
\end{enumerate}

La maquina original, dada su edad, tiene un tiempo medio entre fallas de 57 horas, con un tiempo de reparacion promedio de 19 horas dada su simplicidad. Por otro lado, la nueva maquina posee un tiempo medio entre fallas de 372 horas, muy superior a la maquina antigua. Dada su complejidad, cuando falla es necesario traer un tecnico para repararla, lo cual implica un tiempo medio de reparacion de 124 horas. Considere que todos los tiempos tienen distribucion exponencial y en ambos casos las reparaciones tienen un coeficiente de variabilidad de 1.

\begin{enumerate}[label=\alph*),resume]
  \item Determine la disponibilidad de cada maquina.
  \item Calcule ahora el coeficiente de variabilidad efectivo de cada maquina. Considerando su nueva metrica, cual maquina considera que es mejor y por que?
\end{enumerate}
\end{enunciado}

\begin{pautacompleta}{pauta paso a paso 2024}
Para la maquina original:
\[
  C_{0,1}=\frac{s_1}{t_1}=\frac{2}{12}=0.167.
\]
Para la maquina nueva, primero se convierten 100 segundos a minutos:
\[
  100\ \text{s}=\frac{100}{60}=1.667\ \text{min}.
\]
Entonces:
\[
  C_{0,2}=\frac{1.667}{10}=0.167.
\]
Ambas maquinas tienen el mismo coeficiente de variacion nominal del tiempo de proceso.

La disponibilidad se calcula con:
\[
  A=\frac{MTTF}{MTTF+MTTR}.
\]
Para la maquina original:
\[
  A_1=\frac{57}{57+19}=\frac{57}{76}=0.75.
\]
Para la maquina nueva:
\[
  A_2=\frac{372}{372+124}=\frac{372}{496}=0.75.
\]
Con solo disponibilidad, ambas parecen equivalentes.

Ahora se calcula el coeficiente de variabilidad efectivo:
\[
  C_e^2=C_0^2+(1+C_r^2)A(1-A)\frac{MTTR}{t_0}.
\]
Como las reparaciones tienen coeficiente de variabilidad 1:
\[
  C_r=1.
\]
Para la maquina original:
\[
  C_{e,1}^2=0.1667^2+(1+1^2)(0.75)(1-0.75)\frac{19}{12}.
\]
\[
  C_{e,1}^2=0.0278+2(0.75)(0.25)(1.5833)=0.6215\approx 0.62.
\]
Para la maquina nueva:
\[
  C_{e,2}^2=0.1667^2+(1+1^2)(0.75)(1-0.75)\frac{124}{10}.
\]
\[
  C_{e,2}^2=0.0278+2(0.75)(0.25)(12.4)=4.6778\approx 4.68.
\]
Aunque ambas tienen igual disponibilidad y variabilidad nominal parecida, la maquina nueva genera mucha mas variabilidad efectiva porque sus reparaciones son muy largas en relacion con su tiempo de proceso. Bajo esta metrica, la maquina original es mejor si el objetivo es reducir variabilidad del flujo.
\end{pautacompleta}

\section{Arrastre I1/I2: Recetas que Siguen Apareciendo}

\subsection{Pronosticos}

\begin{enunciado}{Examen 2023, medias moviles y proposito del pronostico}
Un almacen ha decidido sofisticar la forma que hace los pedidos. Para ello la dueña del almacen ha tomado varias decisiones que espera surtan efecto en reducir sus costos y mejorar la atencion a sus clientes. Ahora, usted, como ingeniero industrial experto en gestion de operaciones, debe apoyarla y darle buenas recomendaciones en que hacer. La primera accion que hace es mejorar la forma de gestion de los productos que compran a Carozzi. En particular, quiere gestionar los tallarines grado 1 en paquetes de 400 gramos. A continuacion, esta una tabla con las ventas semanales de las ultimas 5 semanas de este producto.

\begin{center}
\begin{tabular}{c|ccccc}
\toprule
Semana & 1 & 2 & 3 & 4 & 5\\
\midrule
Ventas (unidades) & 36 & 18 & 0 & 20 & 27\\
\bottomrule
\end{tabular}
\end{center}

\begin{enumerate}[label=\alph*)]
  \item Lo primero que ha hecho la dueña es comenzar a usar metodos de estimacion de demanda. En particular, para la semana 6 ha decidido usar medias moviles, pero no sabe si el tamaño de la ventana deberia ser 2 o 3 semanas. Cual de las dos le recomendaria usted?
  \item Como vimos en clase, un pronostico siempre debe venir asociado a un ``para que hacerlo'', es decir, a una finalidad. Conteste:
  \begin{enumerate}[label=\roman*.]
    \item Esta de acuerdo con los datos usados por el dueño del almacen para hacer el pronostico de demanda? Justifique si el pronostico obtenido logra o no el objetivo.
    \item La almacenera no sabe si su pronostico es bueno. Cuales metricas recomendaria utilizar y por que?
  \end{enumerate}
  \item Suponga que la estimacion de demanda para la semana 6 resulta en 25 unidades y la dueña considera que se mantendra constante semanalmente por las proximas 20 semanas. Comprar a Carozzi cada paquete cuesta \$580. Considerando costos de almacenar, costo financiero y mermas, estima que mantener inventario cuesta 7\% del costo del producto por semana. Cada vez que compra, entre tiempo y transporte, gasta \$1.316 por compra y demora 2 semanas en llegar. Cuantas unidades recomienda comprar cada vez para asegurar producto siempre y minimizar costos de gestion de inventario?
\end{enumerate}
\end{enunciado}

\begin{pautacompleta}{pauta paso a paso 2023}
Para comparar ventanas se usa una metrica historica comparable, como MSE o MAPE. La pauta indica que el mejor pronostico historico se obtiene con $n=3$:
\[
  MSE_{n=3}=104.7,\qquad MSE_{n=2}=379.7.
\]
Por lo tanto se recomienda ventana de 3 semanas.

El punto conceptual central es que ventas no siempre es igual a demanda. En la semana 3 hubo venta cero, pero no se sabe si fue por demanda cero o por quiebre de stock. Si no habia inventario, usar esa venta como demanda subestima la demanda futura y contradice el objetivo de no perder ventas.

MSE sirve para medir error promedio, pero no detecta si se esta subestimando sistematicamente. Como el objetivo es no perder ventas, conviene complementar con MFE:
\[
  MFE=\frac{1}{n}\sum_t e_t.
\]
Si $MFE$ es positivo o negativo de forma persistente, muestra sesgo.

Para el inciso c), la demanda total del horizonte es:
\[
  D=25(20)=500.
\]
El costo de mantener una unidad por semana es:
\[
  H=0.07(580)=40.6.
\]
El costo de ordenar es:
\[
  S=1316.
\]
Con EOQ:
\[
  Q^\star=\sqrt{\frac{2DS}{H}}
  =\sqrt{\frac{2(500)(1316)}{40.6}}
  \approx 180.
\]
Se recomienda ordenar aproximadamente 180 unidades por pedido.
\end{pautacompleta}

\subsection{PERT y contratos}

\begin{enunciado}{Examen 2021, bono y penalidad}
Usted se encuentra a cargo de un proyecto y ha determinado que la ruta critica esperada tiene una duracion de 240 dias con una desviacion standard de 20 dias. Actualmente se encuentra cerrando el contrato del proyecto y su contraparte le ofrece un bono de \$300 millones por terminar antes de una fecha y una penalidad de \$200 millones si termina despues de dicha fecha.

\begin{enumerate}[label=\alph*)]
  \item Si la contraparte le ofrece terminar el proyecto en 220 dias, acepta usted el contrato? Muestre todos sus calculos.
  \item Una empresa consultora le asegura que puede reducir el tiempo del proyecto a 225 dias y su desviacion standard a 15 dias a costo fijo de \$10 millones. Utilizando la consultora, aceptaria el contrato de terminar en 220 dias? Muestre todos sus calculos.
\end{enumerate}
\end{enunciado}

\begin{pautacompleta}{pauta paso a paso 2021}
Sin consultora:
\[
  Z=\frac{D-T_E}{\sigma_c}=\frac{220-240}{20}=-1.
\]
\[
  P(T\leq 220)=P(Z\leq -1)=0.1587.
\]
El valor esperado del bono es:
\[
  300(0.1587)=47.61.
\]
El valor esperado de la penalidad es:
\[
  200(1-0.1587)=168.26.
\]
Como $47.61-168.26<0$, no conviene aceptar.

Con consultora:
\[
  Z=\frac{220-225}{15}=-0.33.
\]
\[
  P(T\leq 220)\approx 0.3707.
\]
El valor esperado neto incluyendo costo fijo es:
\[
  VE=300(0.3707)-200(1-0.3707)-10.
\]
\[
  VE=111.21-125.86-10=-24.65.
\]
Tampoco conviene aceptar el contrato bajo estos datos.
\end{pautacompleta}

\begin{enunciado}{Examen 2023, PERT con confianza}
Usted trabaja en una empresa vitivinicola, en donde se han ganado un importante contrato que requiere ampliar su capacidad de produccion de vinos embotellados para el proximo año. Para asegurar la fecha del primer envio al nuevo cliente, usted contrata a una empresa de ingenieria que entrega la secuencia mas logica de construccion y escenarios de duracion.

\begin{center}
\begin{tabular}{c|c|ccc}
\toprule
Actividad & Antecesora & Optimista & Mas probable & Pesimista\\
\midrule
A & - & 1 & 2 & 3\\
B & A & 4 & 5 & 12\\
C & B & 5 & 6 & 7\\
D & B & 1 & 1 & 1\\
E & B,D & 1.5 & 3 & 4.5\\
F & C,E & 2 & 2 & 2\\
\bottomrule
\end{tabular}
\end{center}

\begin{enumerate}[label=\alph*)]
  \item Elabore el diagrama de PERT y determine la ruta critica del proyecto.
  \item Considerando que la construccion parte en un mes mas, y en base a los antecedentes entregados, en cuantos meses mas podria asegurar el primer envio al nuevo cliente usando 95\% de confianza?
  \item Usted va a implementar el sistema del Ultimo Planificador y su cliente quiere estar despachando en la semana 15 con 98\% de confianza. Cual deberia ser la maxima variabilidad permitida para cumplir?
\end{enumerate}
\end{enunciado}

\begin{pautacompleta}{pauta PERT 2023}
Para cada actividad:
\[
  \mu_i=\frac{a_i+4m_i+b_i}{6},\qquad
  \sigma_i=\frac{b_i-a_i}{6},\qquad
  \Var_i=\sigma_i^2.
\]
La red es:
\[
  A\rightarrow B,\qquad B\rightarrow C,\qquad B\rightarrow D,\qquad (B,D)\rightarrow E,\qquad (C,E)\rightarrow F.
\]
La pauta obtiene ruta critica:
\[
  A-B-C-F,
\]
con:
\[
  T_E=16,\qquad \Var_c=2.
\]
Con 95\% de confianza:
\[
  D_{95}=16+1.65\sqrt{2}=18.33\ \text{semanas}.
\]
Si la construccion parte en un mes mas, se suma ese mes al compromiso calendario.

Para semana 15 con 98\%, se exigiria:
\[
  15\geq 16+z_{0.98}\sigma_c.
\]
Como $15<16$, reducir solo la variabilidad no basta: incluso con $\sigma_c=0$, el promedio esperado es 16. Para cumplir semana 15 con mas de 50\% de confianza hay que reducir tambien la duracion esperada.
\end{pautacompleta}

\subsection{Wagner-Whitin y lotificacion}

\begin{enunciado}{Examen 2019, EOQ versus Wagner-Whitin}
La demanda semanal del producto se muestra en la siguiente tabla.

\begin{center}
\begin{tabular}{c|ccccccc|c}
\toprule
Semana & 1 & 2 & 3 & 4 & 5 & 6 & 7 & Promedio\\
\midrule
Pedido & 140 & 90 & 130 & 120 & 100 & 60 & 80 & 102.86\\
\bottomrule
\end{tabular}
\end{center}

Si la maquina que produce el producto tiene un costo de setup de \$100 cada vez que se inicia el proceso y el producto tiene un costo de inventario de \$1 por semana, determine un plan de produccion utilizando EOQ y otro plan utilizando Wagner-Whitin. Cual prefiere y por que se producen las diferencias?
\end{enunciado}

\begin{pautacompleta}{pauta 2019}
Con EOQ se usa demanda promedio:
\[
  D=102.86,\qquad S=100,\qquad H=1.
\]
\[
  Q^\star=\sqrt{\frac{2DS}{H}}\approx 143.33\approx 144.
\]
La pauta reporta un plan EOQ con costo total 950.

Wagner-Whitin usa demanda deterministica por periodo y minimiza setup mas inventario:
\[
  F(t)=\min_{1\le k\le t}\left\{F(k-1)+S+\sum_{j=k}^{t}h(j-k)D_j\right\}.
\]
La pauta reporta un plan WW con costo 669. Se prefiere WW porque produce considerando explicitamente la variabilidad temporal de la demanda; EOQ suaviza por promedio y genera inventario innecesario.
\end{pautacompleta}

\subsection{MRP y modelos de produccion}

\begin{enunciado}{Examen 2017, BOM y programacion matematica}
Usted es el Gerente de Operaciones de una empresa que produce el producto P1. A continuacion se detalla el BOM, donde el numero indica las piezas necesarias:

\begin{itemize}
  \item P1 requiere PO (1) y PU (2).
  \item PO requiere A (3) y B (2).
  \item B requiere C (2).
  \item PU requiere A (1) y C (2).
\end{itemize}

El tiempo de produccion o entrega de cada pieza es $L_k$ semanas para cada pieza $k$, y las maquinas y proveedores tienen una capacidad maxima de entregar $C_k$ unidades a la semana. Los costos semanales de inventario son \$ $I_k$ por cada unidad de pieza $k$ mantenida en inventario y el costo de produccion es \$ $P_k$ por cada unidad de pieza $k$ producida. La demanda por cada unidad de pieza $k$ en el periodo $t$ es $d_{kt}$. Finalmente, las piezas PO y PU comparten la maquina M1 y las piezas A y B comparten la maquina M2, por lo que no pueden ser producidas en la misma semana.

Con esta informacion plantee un problema de programacion matematica que permita obtener el plan de produccion de la empresa.
\end{enunciado}

\begin{pautacompleta}{modelo matematico}
Sea:
\[
  K=\{P1,PO,PU,A,B,C\},\qquad T=\{1,\ldots,H\}.
\]
Variables:
\[
  x_{kt}\geq 0=\text{cantidad de }k\text{ lanzada/producida en }t,
\]
\[
  I_{kt}\geq 0=\text{inventario de }k\text{ al final de }t,
  \qquad
  y_{kt}\in\{0,1\}.
\]
La variable $y_{kt}$ vale 1 si se produce la pieza $k$ en la semana $t$.

La funcion objetivo:
\[
  \min \sum_{k\in K}\sum_{t\in T} P_kx_{kt}
  +\sum_{k\in K}\sum_{t\in T} I_k I_{kt}.
\]
Los balances deben incluir demanda independiente y dependiente:
\[
  I_{kt}=I_{k,t-1}+x_{k,t-L_k}-D_{kt}^{\text{total}},
\]
donde:
\[
  D_{P1,t}^{\text{total}}=d_{P1,t},
\]
\[
  D_{PO,t}^{\text{total}}=d_{PO,t}+x_{P1,t},
  \qquad
  D_{PU,t}^{\text{total}}=d_{PU,t}+2x_{P1,t},
\]
\[
  D_{A,t}^{\text{total}}=d_{A,t}+3x_{PO,t}+x_{PU,t},
\]
\[
  D_{B,t}^{\text{total}}=d_{B,t}+2x_{PO,t},
\]
\[
  D_{C,t}^{\text{total}}=d_{C,t}+2x_{B,t}+2x_{PU,t}.
\]
Capacidad:
\[
  x_{kt}\leq C_k y_{kt}\qquad \forall k,t.
\]
Maquinas compartidas:
\[
  y_{PO,t}+y_{PU,t}\leq 1,\qquad
  y_{A,t}+y_{B,t}\leq 1.
\]
No negatividad e integralidad:
\[
  x_{kt},I_{kt}\geq 0,\qquad y_{kt}\in\{0,1\}.
\]
\end{pautacompleta}

\subsection{Localizacion y punto de equilibrio}

\begin{enunciado}{Examen 2024, localizacion L1 versus L2}
Usted se encuentra seleccionando la localizacion para su planta de produccion. Actualmente tiene dos lugares como posibles candidatos L1 y L2. Los valores de compra del terreno y costos variables de operacion para cada localizacion son:

\begin{center}
\begin{tabular}{c|cc}
\toprule
Localizacion & Costo Compra (U\$) & Costo Operacion (U\$/unidad)\\
\midrule
L1 & 10.000 & 25\\
L2 & 40.000 & 10\\
\bottomrule
\end{tabular}
\end{center}

\begin{enumerate}[label=\alph*)]
  \item Que nivel de produccion lo deja indiferente entre L1 o L2?
  \item Usted sabe que hoy, tiempo 0, tiene demanda de 1.000 unidades, estima que en tiempo 9 tendra demanda de 3.100 unidades y la demanda crece linealmente. Si puede colocarse en L1 y despues cambiarse a L2, pero a un costo fijo de cambio, que maximo costo fijo estaria dispuesto a aceptar?
  \item Si ahora solo puede decidir por una ubicacion y la demanda es $Q(t)=1000+250t$, sin cambio de valor del dinero, desde $t=0$ a $t=10$, que ubicacion es mas adecuada?
\end{enumerate}
\end{enunciado}

\begin{pautacompleta}{pauta 2024}
Punto de indiferencia:
\[
  10000+25X=40000+10X
  \Rightarrow X=2000.
\]
Para el cambio, el beneficio aparece despues de cruzar el punto de equilibrio. La pauta calcula:
\[
  \text{Lugar 1}=\frac{(3100-2000)(87500-60000)}{2}=15{,}125{,}000,
\]
\[
  \text{Lugar 2}=\frac{(3100-2000)(71000-60000)}{2}=6{,}050{,}000.
\]
El diferencial es:
\[
  15{,}125{,}000-6{,}050{,}000=9{,}075{,}000.
\]
Ese es el maximo costo fijo de cambio aceptable.

Para elegir una sola ubicacion:
\[
  C_1(t)=10000+25(1000+250t)=35000+6250t,
\]
\[
  C_2(t)=40000+10(1000+250t)=50000+2500t.
\]
Integrando de 0 a 10:
\[
  \int_0^{10}C_1(t)dt=662500,
  \qquad
  \int_0^{10}C_2(t)dt=625000.
\]
Conviene L2.
\end{pautacompleta}

\subsection{Inventario EOQ con entrega gradual}

\begin{enunciado}{Examen 2016, insumos de autos de juguete}
La empresa Jot Wilz fabrica autos de juguete. Cada auto requiere una unidad de aluminio y 4 ruedas. El precio de una unidad de aluminio es \$500 y cada rueda cuesta \$100. El proveedor demora 2 dias habiles en entregar pedidos. Cada orden cuesta \$50.000 y mantener inventario a la semana cuesta 10\% del precio de cada insumo. Con demanda diaria de 2.400 autos:

\begin{enumerate}[label=\alph*)]
  \item Asumiendo pedidos instantaneos, calcule las cantidades optimas a pedir de aluminio y ruedas. Cual es el punto de reorden de cada insumo?
  \item Calcule el costo anual total asociado.
  \item Si el proveedor ya no entrega pedidos completos sino a tasa de 4.000 unidades/dia para aluminio y 12.000 ruedas/dia, calcule cuanto pedir ahora y cuanto cuesta el cambio.
\end{enumerate}
\end{enunciado}

\begin{pautacompleta}{pauta 2016}
Demanda anual de autos:
\[
  D=2400(5)(52)=624000.
\]
Para aluminio:
\[
  H_A=500(0.1)(52).
\]
\[
  Q_A^\star=\sqrt{\frac{2(624000)(50000)}{500(0.1)(52)}}=4898.98\approx 4899.
\]
Para ruedas:
\[
  D_R=4(624000),\qquad H_R=100(0.1)(52).
\]
\[
  Q_R^\star=\sqrt{\frac{2(624000)(4)(50000)}{100(0.1)(52)}}=21908.9\approx 21909.
\]
Puntos de reorden:
\[
  R_A=2(2400)=4800,
  \qquad
  R_R=2(2400)(4)=19200.
\]
Costos anuales:
\[
  CT_A=624000(500)+\frac{624000(50000)}{4899}
  +\frac{4899(500)(0.1)(52)}{2}=324{,}737{,}347.
\]
\[
  CT_R=624000(4)(100)+\frac{624000(4)(50000)}{21909}
  +\frac{21909(100)(0.1)(52)}{2}=260{,}992{,}629.
\]
\[
  CT_{\text{total}}=585{,}729{,}976.
\]
Con entrega gradual:
\[
  Q^\star_{\text{gradual}}=\sqrt{\frac{2DS}{H}}\sqrt{\frac{p}{p-d}}.
\]
Para aluminio la pauta reporta:
\[
  Q_A=7746.
\]
La razon conceptual es que al recibir gradualmente y consumir simultaneamente, el inventario maximo baja respecto de $Q$, y eso cambia el lote economico.
\end{pautacompleta}

\subsection{Newsvendor y productos perecibles}

\begin{enunciado}{Examen 2017, parkas por temporada}
Usted tiene una tienda de ropa mayorista y esta preocupado porque no puso atencion en clases cuando se vieron los modelos de inventarios. Tiene un problema con un producto nuevo: parkas. La venta se hace en lotes de 25 y marketing entrega la distribucion de probabilidades de venta:

\begin{center}
\begin{tabular}{c|cccccc}
\toprule
Cantidad parkas & 300 & 325 & 350 & 375 & 400 & 425\\
\midrule
Probabilidad & 0.1 & 0.2 & 0.3 & 0.24 & 0.1 & 0.06\\
\bottomrule
\end{tabular}
\end{center}

Cada parka se vende en temporada a \$130, el proveedor cobra \$100 y si no se vende debe liquidarse en outlet a \$80.

\begin{enumerate}[label=\alph*)]
  \item Que cantidad de parkas deberia pedir? Cual es la ganancia/perdida esperada?
  \item Suponiendo que la distribucion subyacente es normal, cambia la cantidad a pedir?
\end{enumerate}
\end{enunciado}

\begin{pautacompleta}{pauta newsvendor}
Costo de quedarse corto:
\[
  C_u=130-100=30.
\]
Costo de quedarse largo:
\[
  C_o=100-80=20.
\]
Fractil critico:
\[
  \frac{C_u}{C_u+C_o}=\frac{30}{50}=0.6.
\]
Distribucion acumulada:
\[
  F(300)=0.1,\quad F(325)=0.3,\quad F(350)=0.6.
\]
El menor lote con acumulada al menos 0.6 es:
\[
  Q^\star=350.
\]
La ganancia esperada:
\[
  \E[\Pi(350)]=
  \sum_d p_d\left[130\min(350,d)+80(350-d)^+-100(350)\right].
\]
Si se asume normal, se calcula:
\[
  \mu=\sum d p_d,\qquad \sigma^2=\sum(d-\mu)^2p_d,
\]
y:
\[
  Q_N^\star=\mu+z_{0.6}\sigma.
\]
Luego se redondea al multiplo de 25 mas cercano que maximice utilidad esperada.
\end{pautacompleta}

\begin{enunciado}{Examen 2022, producto perecible premium/secundario}
Usted esta a cargo del proceso productivo que produce un producto $P$ que usa solo un insumo $I$. Para producir una unidad de $P$ se requiere una unidad de $I$. El proceso se compone de dos maquinas en serie. La maquina A toma el insumo $I$ y lo procesa a una tasa de 1 segundo por unidad, produciendo un producto semiterminado. Este se almacena en una camara de temperatura controlada para ser procesado por la maquina B a 50 unidades por minuto. El proceso opera 24 horas al dia los 365 dias del año.

El producto final se vende al mercado premium a \$150 por unidad. Si no es vendido durante el dia, debe liquidarse en mercado secundario a \$80 por unidad. La demanda es normal con media 68.000 unidades/dia y desviacion standard 12.000. El costo total del producto terminado es \$100 por unidad. El costo de almacenamiento anual es \$17 por unidad por dia para producto final, \$15 por unidad al año para semiterminado y \$7.3 por unidad para insumo $I$. Emitir una orden al proveedor cuesta \$200.000 y el proveedor demora 5 dias en entregar.

\begin{enumerate}[label=\roman*.]
  \item Determine la cantidad diaria a producir y el inventario de producto terminado.
  \item Cual es el lote optimo de pedido del insumo $I$ y su punto de reorden?
  \item Si solo la maquina A tiene costo de setup $CT_{\text{Setup}}=300000/Q$, desarrolle una forma funcional del costo total de pedido del semiterminado y plantee la ecuacion del lote optimo.
\end{enumerate}
\end{enunciado}

\begin{pautacompleta}{pauta 2022}
Newsvendor:
\[
  C_u=150-100=50,\qquad C_o=100-80=20.
\]
\[
  F(Q^\star)=\frac{50}{50+20}=0.7142.
\]
La pauta usa $z=0.57$:
\[
  Q^\star=68000+0.57(12000)=74840.
\]
Capacidad:
\[
  A=24(60)(60)=86400\ \text{unid/dia},
\]
\[
  B=50(60)(24)=72000\ \text{unid/dia}.
\]
El cuello es B, por lo que:
\[
  Q=72000\ \text{unid/dia}.
\]
La pauta deja inventario terminado en 0 porque se produce y vende diariamente.

Para insumo:
\[
  R=dL=72000(5)=360000.
\]
El lote optimo se calcula con EOQ:
\[
  Q^\star=\sqrt{\frac{2DS}{H}}.
\]
Para semiterminado:
\[
  CT(Q)=DC+\frac{300000D}{Q^2}
  +\frac{1}{2}\left(1-\frac{d}{p}\right)QH.
\]
La condicion de optimalidad:
\[
  -2\frac{300000D}{Q^3}
  +\frac{1}{2}\left(1-\frac{d}{p}\right)H=0.
\]
\end{pautacompleta}

\subsection{Centralizacion de inventarios}

\begin{enunciado}{Examen 2017, dos mercados y centralizacion}
Usted esta a cargo de la logistica de una empresa que vende al retail. Actualmente la empresa cuenta con un producto y dos mercados. El costo de cada orden es \$60 y mantener inventario cuesta \$0.27 por unidad a la semana. El gerente pide nivel de servicio 97\% ($z=1.88$). La fabrica tiene lead time de 1 semana. La demanda es:

\begin{center}
\begin{tabular}{c|ccccc|cc}
\toprule
Mercado & Semana 1 & 2 & 3 & 4 & 5 & Promedio & Desv. Est.\\
\midrule
1 & 33 & 45 & 37 & 38 & 55 & 41.6 & 8.6\\
2 & 46 & 35 & 41 & 40 & 26 & 37.6 & 7.6\\
\bottomrule
\end{tabular}
\end{center}

Se despacha a cada mercado independientemente y se analiza centralizar. El transporte descentralizado cuesta \$1.05 por unidad y centralizado \$1.10 por unidad; el precio de venta es \$1.

\begin{enumerate}[label=\alph*)]
  \item Determine politica descentralizada para cada mercado bajo revision continua y costo anual.
  \item Determine politica centralizada bajo revision continua y costo anual.
  \item Recomiende centralizar o mantener descentralizado. Para $N$ productos y $M$ mercados, construya un modelo matematico para decidir.
\end{enumerate}
\end{enunciado}

\begin{pautacompleta}{pauta pooling}
Para cada mercado:
\[
  Q_m^\star=\sqrt{\frac{2D_mS}{H}},\qquad
  R_m=d_mL+z\sigma_m\sqrt{L}.
\]
Con semanas como unidad:
\[
  R_1=41.6+1.88(8.6),\qquad
  R_2=37.6+1.88(7.6).
\]
Centralizado:
\[
  d_C=41.6+37.6=79.2.
\]
Si las demandas son independientes:
\[
  \sigma_C=\sqrt{8.6^2+7.6^2}=11.48.
\]
\[
  R_C=79.2+1.88(11.48).
\]
El beneficio potencial de centralizar viene del risk pooling: la desviacion agregada no es $8.6+7.6$, sino la raiz de la suma de varianzas. La decision compara costo de inventario, ordenes y transporte:
\[
  CT_D \quad \text{versus}\quad CT_C.
\]
Para $N$ productos y $M$ mercados, se puede definir $y_i=1$ si producto $i$ se centraliza y minimizar:
\[
  \min \sum_i \left[y_i CT_{i,C}+(1-y_i)CT_{i,D}\right],
\]
con $y_i\in\{0,1\}$ y restricciones de capacidad si existen.
\end{pautacompleta}

\subsection{Regresion de demanda y lotes con maquinas en serie}

\begin{enunciado}{Examen 2018, acido sulfurico}
Usted es gerente de una planta de acido sulfurico que provee a la mineria. La demanda diaria de acido en toneladas depende del precio diario del cobre $PCu$ en U\$ por libra.

\begin{center}
\begin{tabular}{c|cc|c|cc}
\toprule
\# & Precio Cu & Q acido & \# & Precio Cu & Q acido\\
\midrule
1&3.2&1200&10&3.5&1450\\
2&3.0&1000&11&2.6&750\\
3&3.1&1050&12&3.1&1100\\
4&3.3&1270&13&3.2&1200\\
5&2.7&800&14&2.9&990\\
6&3.4&1350&15&2.6&750\\
7&3.0&1025&16&2.8&900\\
8&2.8&850&17&3.0&1025\\
9&2.9&950&18&3.1&1050\\
\bottomrule
\end{tabular}
\end{center}

Se tiene:
\[
  \sum PCu=54.2,\quad \sum Q=18710,\quad \sum PCu^2=164.3,
\]
\[
  \sum Q^2=20114250,\quad \sum Q\cdot PCu=57192.
\]
El proceso tiene dos maquinas en serie:
\[
  M1:\;0.016\ \text{hr/ton},\qquad M2:\;40\ \text{ton/hr}.
\]
Opera 3 turnos, 24 horas al dia, 7 dias a la semana. Cada vez que se enciende una maquina hay setup: M1 cuesta \$80/setup y M2 \$120/setup. Las maquinas no ajustan capacidad y funcionan siempre al maximo. Mantener inventario en proceso cuesta \$30 por tonelada al año.

\begin{enumerate}[label=\alph*)]
  \item Determine politica optima de inventarios en proceso/terminado y tiempos de produccion para precio promedio 3.0 U\$/lb de cobre.
  \item Para escenario optimista 3.3 U\$/lb y pesimista 2.7 U\$/lb, calcule nuevamente la politica.
\end{enumerate}
\end{enunciado}

\begin{pautacompleta}{receta regresion + lote economico}
Ajuste:
\[
  \hat Q=a+bP.
\]
\[
  b=\frac{n\sum PQ-\sum P\sum Q}{n\sum P^2-(\sum P)^2},
  \qquad
  a=\bar Q-b\bar P.
\]
Con $n=18$ se pronostica demanda para $P=3.0$, $3.3$ y $2.7$.

Capacidades:
\[
  M1:\frac{1}{0.016}=62.5\ \text{ton/hr},\qquad M2=40\ \text{ton/hr}.
\]
M2 es el cuello. Como las maquinas producen a maxima capacidad al encenderse, se usa lote economico con tasa finita:
\[
  Q_j^\star=\sqrt{\frac{2dS_j}{H}}\sqrt{\frac{p_j}{p_j-d}},
\]
si $p_j>d$. El tiempo de corrida:
\[
  t_j=\frac{Q_j^\star}{p_j}.
\]
La receta es: estimar $d$ con regresion para cada escenario, verificar que $d<p_j$, calcular lote por maquina con setup e inventario, y reportar tiempos de produccion e inventario esperado.
\end{pautacompleta}


\end{document}
